\documentclass{beamer}
\usetheme{Madrid}
\usecolortheme{default}
\usepackage[utf8]{inputenc}
\usepackage{amsmath}
\usepackage{amsfonts}
\usepackage{amssymb}
\usepackage{tikz}

\title{CRISC Hard Questions from ISACA QAE Manual}
\subtitle{Expanded Bank --- Answers on the Next Page}
\author{Compiled from ISACA CRISC Review Questions, Answers \& Explanations Manual, 6th Edition}
\date{\today}

\begin{document}

	% ============================================================
% DOMAIN 1 --- MORE GOVERNANCE QUESTIONS (Q416–Q440)
% ============================================================
\section{Domain 1: More Governance}

\begin{frame}{D1-Q416: Risk Appetite vs Risk Tolerance}
	\begin{block}{Question}
		Which statement BEST distinguishes risk appetite from risk tolerance?
	\end{block}
	\begin{enumerate}[(A)]
		\item Appetite is a measurable limit; tolerance is a broad statement
		\item Appetite is the amount of risk the enterprise is willing to accept; tolerance is the acceptable variation around that appetite
		\item Appetite applies only to financial risk; tolerance applies to operational risk
		\item Appetite is set by the CRO; tolerance is set by regulators
	\end{enumerate}
\end{frame}
\begin{frame}{D1-Q416: Answer}
	\begin{block}{Correct Answer: B}
		Appetite is the broad willingness to accept risk in pursuit of objectives. Tolerance translates appetite into specific, measurable limits of acceptable variation for individual risk categories.
	\end{block}
\end{frame}

\begin{frame}{D1-Q417: Risk Capacity vs Risk Appetite}
	\begin{block}{Question}
		Which concept represents the maximum risk an enterprise can absorb before its existence is threatened?
	\end{block}
	\begin{enumerate}[(A)]
		\item Risk appetite
		\item Risk tolerance
		\item Risk capacity
		\item Risk acceptance
	\end{enumerate}
\end{frame}
\begin{frame}{D1-Q417: Answer}
	\begin{block}{Correct Answer: C --- Risk capacity}
		Risk capacity is the maximum risk the enterprise can survive. It is always greater than or equal to appetite; breaching capacity threatens viability.
	\end{block}
\end{frame}

\begin{frame}{D1-Q418: Board Accountability for Risk}
	\begin{block}{Question}
		Who is ultimately accountable to shareholders for enterprise risk management?
	\end{block}
	\begin{enumerate}[(A)]
		\item Chief risk officer
		\item Board of directors
		\item Internal audit
		\item Risk management department
	\end{enumerate}
\end{frame}
\begin{frame}{D1-Q418: Answer}
	\begin{block}{Correct Answer: B}
		The board of directors is ultimately accountable to shareholders, customers, employees, and the public for enterprise risk management.
	\end{block}
\end{frame}

\begin{frame}{D1-Q419: Risk Appetite Communication}
	\begin{block}{Question}
		What is the MOST effective way to communicate risk appetite to operational staff?
	\end{block}
	\begin{enumerate}[(A)]
		\item Email the full appetite statement once a year
		\item Translate appetite into operational tolerances, thresholds, and decision rules
		\item Discuss it only in executive meetings
		\item Post it on the intranet without explanation
	\end{enumerate}
\end{frame}
\begin{frame}{D1-Q419: Answer}
	\begin{block}{Correct Answer: B}
		Operational staff need concrete tolerances, thresholds, and decision rules derived from appetite so they can make risk-aware decisions in daily work.
	\end{block}
\end{frame}

\begin{frame}{D1-Q420: Three Lines of Defense - First Line}
	\begin{block}{Question}
		Which activity is a FIRST line of defense responsibility?
	\end{block}
	\begin{enumerate}[(A)]
		\item Independent assurance on internal controls
		\item Operational risk management and control implementation
		\item Establishing the enterprise risk framework
		\item Board oversight of risk appetite
	\end{enumerate}
\end{frame}
\begin{frame}{D1-Q420: Answer}
	\begin{block}{Correct Answer: B}
		The first line (operational management) owns risk and implements controls. The second line establishes the framework; the third line provides independent assurance.
	\end{block}
\end{frame}

\begin{frame}{D1-Q421: Three Lines of Defense - Second Line}
	\begin{block}{Question}
		Which is a SECOND line of defense activity?
	\end{block}
	\begin{enumerate}[(A)]
		\item Performing day-to-day controls
		\item Risk oversight, framework development, and compliance monitoring
		\item Internal audit testing
		\item Board risk committee meetings
	\end{enumerate}
\end{frame}
\begin{frame}{D1-Q421: Answer}
	\begin{block}{Correct Answer: B}
		The second line provides risk oversight, develops frameworks, monitors compliance, and challenges the first line. Internal audit is the third line.
	\end{block}
\end{frame}

\begin{frame}{D1-Q422: Three Lines of Defense - Third Line}
	\begin{block}{Question}
		What is the PRIMARY role of the third line of defense?
	\end{block}
	\begin{enumerate}[(A)]
		\item Implement controls
		\item Establish risk appetite
		\item Provide independent assurance on governance, risk, and control
		\item Approve risk treatment plans
	\end{enumerate}
\end{frame}
\begin{frame}{D1-Q422: Answer}
	\begin{block}{Correct Answer: C}
		The third line—internal audit—provides independent, objective assurance on the effectiveness of governance, risk management, and internal controls.
	\end{block}
\end{frame}

\begin{frame}{D1-Q423: Risk Governance Failures}
	\begin{block}{Question}
		Which is the MOST common cause of risk governance failures?
	\end{block}
	\begin{enumerate}[(A)]
		\item Insufficient risk management software
		\item Unclear accountability and weak oversight
		\item Too many risk assessments performed
		\item Excessive regulatory guidance
	\end{enumerate}
\end{frame}
\begin{frame}{D1-Q423: Answer}
	\begin{block}{Correct Answer: B}
		Governance failures typically stem from unclear accountability, weak oversight, and poorly defined roles—not tooling or assessment frequency.
	\end{block}
\end{frame}

\begin{frame}{D1-Q424: Blame Culture Inhibitor}
	\begin{block}{Question}
		Which is the MOST effective inhibitor of relevant risk communication in an enterprise?
	\end{block}
	\begin{enumerate}[(A)]
		\item Lack of advanced tools
		\item A blame culture that discourages reporting
		\item Too many risk reports
		\item Excessive regulation
	\end{enumerate}
\end{frame}
\begin{frame}{D1-Q424: Answer}
	\begin{block}{Correct Answer: B}
		A blame culture suppresses open reporting and escalation, making it the single most damaging inhibitor of effective risk communication.
	\end{block}
\end{frame}

\begin{frame}{D1-Q425: Risk Culture Symptom}
	\begin{block}{Question}
		Which is a symptom of a problematic risk culture?
	\end{block}
	\begin{enumerate}[(A)]
		\item Open communication about risk
		\item Overconfidence in the ability to manage risk
		\item Clear understanding of risk appetite
		\item Regular risk reporting to management
	\end{enumerate}
\end{frame}
\begin{frame}{D1-Q425: Answer}
	\begin{block}{Correct Answer: B}
		Overconfidence, dismissal of risk findings, and resistance to challenge are classic symptoms of a weak risk culture.
	\end{block}
\end{frame}

\begin{frame}{D1-Q426: Risk Policy Hierarchy}
	\begin{block}{Question}
		Which document BEST expresses the high-level intent of an organization's security practices?
	\end{block}
	\begin{enumerate}[(A)]
		\item Procedures
		\item Policies
		\item Guidelines
		\item Standards
	\end{enumerate}
\end{frame}
\begin{frame}{D1-Q426: Answer}
	\begin{block}{Correct Answer: B --- Policies}
		Policies express management intent at a high level. Procedures are detailed steps; standards are measurable requirements; guidelines are recommendations.
	\end{block}
\end{frame}

\begin{frame}{D1-Q427: Risk Awareness Primary Focus}
	\begin{block}{Question}
		What is the PRIMARY consideration when developing an IT risk awareness program?
	\end{block}
	\begin{enumerate}[(A)]
		\item Why IT owns technology risk
		\item How technology risk impacts each attendee's business area
		\item How to transfer technology risk
		\item Why technology risk is harder than other risks
	\end{enumerate}
\end{frame}
\begin{frame}{D1-Q427: Answer}
	\begin{block}{Correct Answer: B}
		Stakeholders must understand how IT-related risk affects their specific business area and objectives. IT does not own technology risk; the business does.
	\end{block}
\end{frame}

\begin{frame}{D1-Q428: Governance Accountability Principle}
	\begin{block}{Question}
		Which principle is MOST fundamental to good governance?
	\end{block}
	\begin{enumerate}[(A)]
		\item Centralizing all decisions with the CIO
		\item Clear accountability aligned with authority
		\item Outsourcing all risk decisions
		\item Eliminating all risk
	\end{enumerate}
\end{frame}
\begin{frame}{D1-Q428: Answer}
	\begin{block}{Correct Answer: B}
		Accountability must align with authority. Without this alignment, accountability is unenforceable and governance breaks down.
	\end{block}
\end{frame}

\begin{frame}{D1-Q429: Risk Register Ownership}
	\begin{block}{Question}
		Who should own the enterprise risk register?
	\end{block}
	\begin{enumerate}[(A)]
		\item The internal audit team
		\item The risk management function, with input from risk owners
		\item The IT department alone
		\item An external consultant
	\end{enumerate}
\end{frame}
\begin{frame}{D1-Q429: Answer}
	\begin{block}{Correct Answer: B}
		The risk management function typically maintains the register, with risk owners providing and validating entries for their areas.
	\end{block}
\end{frame}

\begin{frame}{D1-Q430: Risk Appetite Revision}
	\begin{block}{Question}
		Who should approve changes to the enterprise risk appetite?
	\end{block}
	\begin{enumerate}[(A)]
		\item Risk management department
		\item Board of directors with senior management input
		\item Individual business unit managers
		\item Internal audit
	\end{enumerate}
\end{frame}
\begin{frame}{D1-Q430: Answer}
	\begin{block}{Correct Answer: B}
		Risk appetite is a board-level decision informed by senior management, reflecting the enterprise's strategic direction and tolerance.
	\end{block}
\end{frame}

\begin{frame}{D1-Q431: Risk Governance and Strategy}
	\begin{block}{Question}
		How should risk governance relate to business strategy?
	\end{block}
	\begin{enumerate}[(A)]
		\item Risk governance should be independent and unrelated
		\item Risk governance should enable strategy by managing risks to objectives
		\item Risk governance should block strategic change
		\item Risk governance should replace strategic planning
	\end{enumerate}
\end{frame}
\begin{frame}{D1-Q431: Answer}
	\begin{block}{Correct Answer: B}
		Risk governance enables strategy by ensuring risks to objectives are identified, assessed, and managed within appetite—not by blocking change.
	\end{block}
\end{frame}

\begin{frame}{D1-Q432: Board Risk Reporting}
	\begin{block}{Question}
		Which characteristic makes risk reporting MOST useful to the board?
	\end{block}
	\begin{enumerate}[(A)]
		\item Inclusion of every identified risk
		\item Aggregated exposures against appetite with clear trends
		\item Use of technical jargon
		\item Weekly frequency
	\end{enumerate}
\end{frame}
\begin{frame}{D1-Q432: Answer}
	\begin{block}{Correct Answer: B}
		Boards need aggregated, trend-oriented views against appetite—not exhaustive detail—to provide strategic oversight.
	\end{block}
\end{frame}

\begin{frame}{D1-Q433: Risk Culture and Incentives}
	\begin{block}{Question}
		How do incentive structures influence risk culture?
	\end{block}
	\begin{enumerate}[(A)]
		\item They have no effect
		\item They align behavior with desired risk-taking boundaries
		\item They eliminate the need for policies
		\item They prevent all incidents
	\end{enumerate}
\end{frame}
\begin{frame}{D1-Q433: Answer}
	\begin{block}{Correct Answer: B}
		Well-aligned incentives reward appropriate risk-taking and discourage recklessness. Misaligned incentives undermine even the best controls.
	\end{block}
\end{frame}

\begin{frame}{D1-Q434: Risk Governance Maturity}
	\begin{block}{Question}
		Which characteristic indicates a mature risk governance structure?
	\end{block}
	\begin{enumerate}[(A)]
		\item Risk is managed only after incidents
		\item Risk is integrated into strategic planning and decision-making
		\item Risk is managed exclusively by the risk department
		\item Risk reports are annual compliance artifacts
	\end{enumerate}
\end{frame}
\begin{frame}{D1-Q434: Answer}
	\begin{block}{Correct Answer: B}
		Mature governance integrates risk into strategy, budgeting, and operations. Reactive or siloed risk management indicates lower maturity.
	\end{block}
\end{frame}

\begin{frame}{D1-Q435: Third-Party Governance}
	\begin{block}{Question}
		Which is MOST important for enterprise-level third-party risk governance?
	\end{block}
	\begin{enumerate}[(A)]
		\item Delegating oversight to the vendor
		\item A centralized third-party risk program with clear accountability
		\item Annual contract reviews only
		\item Relying on vendor self-assessments
	\end{enumerate}
\end{frame}
\begin{frame}{D1-Q435: Answer}
	\begin{block}{Correct Answer: B}
		A centralized program ensures consistent due diligence, contracting, monitoring, and termination practices across all third parties.
	\end{block}
\end{frame}

\begin{frame}{D1-Q436: Risk Governance and Risk Aggregation}
	\begin{block}{Question}
		Why is risk aggregation important for enterprise governance?
	\end{block}
	\begin{enumerate}[(A)]
		\item It eliminates the need for business-unit risk registers
		\item It provides a holistic view of exposures that may exceed appetite
		\item It ensures every risk has a unique owner
		\item It replaces risk assessment
	\end{enumerate}
\end{frame}
\begin{frame}{D1-Q436: Answer}
	\begin{block}{Correct Answer: B}
		Aggregation reveals concentrations and cumulative exposures that may individually appear acceptable but collectively exceed appetite—critical for board oversight.
	\end{block}
\end{frame}

\begin{frame}{D1-Q437: Governance and Regulatory Integration}
	\begin{block}{Question}
		What is the BEST way to integrate regulatory requirements into risk governance?
	\end{block}
	\begin{enumerate}[(A)]
		\item Treat compliance as separate from risk management
		\item Map regulatory obligations to risk categories and controls within the ERM framework
		\item Wait for regulators to identify gaps
		\item Delegate compliance entirely to external counsel
	\end{enumerate}
\end{frame}
\begin{frame}{D1-Q437: Answer}
	\begin{block}{Correct Answer: B}
		Mapping obligations to risk categories and controls ensures compliance is managed proactively within the ERM framework—not as an isolated activity.
	\end{block}
\end{frame}

\begin{frame}{D1-Q438: Governance Effectiveness Measure}
	\begin{block}{Question}
		Which is the BEST measure of risk governance effectiveness?
	\end{block}
	\begin{enumerate}[(A)]
		\item Number of policies published
		\item Degree to which risk-informed decisions are made and evidenced
		\item Size of the risk management department
		\item Frequency of risk committee meetings
	\end{enumerate}
\end{frame}
\begin{frame}{D1-Q438: Answer}
	\begin{block}{Correct Answer: B}
		Effectiveness is measured by whether risk information actually informs decisions—not by documentation volume, headcount, or meeting frequency.
	\end{block}
\end{frame}

\begin{frame}{D1-Q439: Risk Reporting Frequency}
	\begin{block}{Question}
		What is the BEST approach for determining risk reporting frequency to senior management?
	\end{block}
	\begin{enumerate}[(A)]
		\item Always monthly
		\item Aligned with the pace of change and criticality of reported risks
		\item Only annually
		\item Only after incidents
	\end{enumerate}
\end{frame}
\begin{frame}{D1-Q439: Answer}
	\begin{block}{Correct Answer: B}
		Frequency should reflect how rapidly the risk environment changes and how critical those risks are to business objectives.
	\end{block}
\end{frame}

\begin{frame}{D1-Q440: Governance and Emerging Risk}
	\begin{block}{Question}
		How should emerging risk be governed at the enterprise level?
	\end{block}
	\begin{enumerate}[(A)]
		\item Ignore until it becomes a realized loss
		\item Establish a process for identifying, assessing, and escalating emerging risk
		\item Delegate entirely to IT
		\item Report only to external auditors
	\end{enumerate}
\end{frame}
\begin{frame}{D1-Q440: Answer}
	\begin{block}{Correct Answer: B}
		A structured process for identifying, assessing, and escalating emerging risk enables proactive management before issues crystallize into material losses.
	\end{block}
\end{frame}

% ============================================================
% DOMAIN 2 --- MORE RISK ASSESSMENT (Q441–Q465)
% ============================================================
\section{Domain 2: More Risk Assessment}

\begin{frame}{D2-Q441: Risk Assessment Objectives}
	\begin{block}{Question}
		What is the PRIMARY objective of a risk assessment?
	\end{block}
	\begin{enumerate}[(A)]
		\item To document every possible risk
		\item To identify and evaluate risk to enable risk-aware decisions
		\item To eliminate all risk
		\item To satisfy auditors
	\end{enumerate}
\end{frame}
\begin{frame}{D2-Q441: Answer}
	\begin{block}{Correct Answer: B}
		Risk assessment identifies and evaluates risk and its potential effects to enable informed, risk-aware decision-making.
	\end{block}
\end{frame}

\begin{frame}{D2-Q442: Asset Identification First}
	\begin{block}{Question}
		What is the FIRST step in the risk assessment process?
	\end{block}
	\begin{enumerate}[(A)]
		\item Identify threats
		\item Identify assets
		\item Identify vulnerabilities
		\item Identify threat sources
	\end{enumerate}
\end{frame}
\begin{frame}{D2-Q442: Answer}
	\begin{block}{Correct Answer: B}
		Asset identification comes first—you cannot assess risk to what you have not identified. Threats and vulnerabilities follow.
	\end{block}
\end{frame}

\begin{frame}{D2-Q443: Risk Scenario Analysis Purpose}
	\begin{block}{Question}
		Risk scenarios are analyzed to determine:
	\end{block}
	\begin{enumerate}[(A)]
		\item Strength of controls
		\item Likelihood and impact
		\item Current risk profile
		\item Root cause
	\end{enumerate}
\end{frame}
\begin{frame}{D2-Q443: Answer}
	\begin{block}{Correct Answer: B}
		Risk scenarios describe events that can lead to business impact and are evaluated for likelihood and impact should the event occur.
	\end{block}
\end{frame}

\begin{frame}{D2-Q444: Risk Register Content}
	\begin{block}{Question}
		Which information is ESSENTIAL for each risk in a risk register?
	\end{block}
	\begin{enumerate}[(A)]
		\item Date, description, impact, probability, score, mitigation, owner
		\item Date, description, cost, communication plan, owner
		\item Description, cost, owner
		\item Risk, controls, audit trail
	\end{enumerate}
\end{frame}
\begin{frame}{D2-Q444: Answer}
	\begin{block}{Correct Answer: A}
		A complete register entry includes date, description, impact, probability, risk score, mitigation action, and owner.
	\end{block}
\end{frame}

\begin{frame}{D2-Q445: Quantitative Risk Analysis Drawback}
	\begin{block}{Question}
		The MOST significant drawback of quantitative risk analysis vs. qualitative is:
	\end{block}
	\begin{enumerate}[(A)]
		\item Lower objectivity
		\item Greater reliance on expertise
		\item Less management buy-in
		\item Higher cost
	\end{enumerate}
\end{frame}
\begin{frame}{D2-Q445: Answer}
	\begin{block}{Correct Answer: D}
		Quantitative analysis is more complex and costly, requiring reliable data and statistical effort.
	\end{block}
\end{frame}

\begin{frame}{D2-Q446: Statistical Methods in Risk Analysis}
	\begin{block}{Question}
		Deriving likelihood and impact through statistical methods is MOST associated with:
	\end{block}
	\begin{enumerate}[(A)]
		\item Risk scenario
		\item Qualitative
		\item Quantitative
		\item Semiquantitative
	\end{enumerate}
\end{frame}
\begin{frame}{D2-Q446: Answer}
	\begin{block}{Correct Answer: C}
		Quantitative risk analysis derives likelihood and impact from statistical methods and data.
	\end{block}
\end{frame}

\begin{frame}{D2-Q447: BIA Input to BCP}
	\begin{block}{Question}
		Which is a PRIMARY input to the Business Continuity Plan?
	\end{block}
	\begin{enumerate}[(A)]
		\item RTO
		\item DR test results
		\item Business Impact Analysis (BIA)
		\item BC test results
	\end{enumerate}
\end{frame}
\begin{frame}{D2-Q447: Answer}
	\begin{block}{Correct Answer: C}
		The BIA identifies critical processes, impacts, and recovery priorities; these outputs are the foundation for the BCP.
	\end{block}
\end{frame}

\begin{frame}{D2-Q448: RPO Meaning}
	\begin{block}{Question}
		A manager confirms an application cannot lose more than four hours of data. Which objective is this?
	\end{block}
	\begin{enumerate}[(A)]
		\item MTD
		\item RPO
		\item RTO
		\item BIA
	\end{enumerate}
\end{frame}
\begin{frame}{D2-Q448: Answer}
	\begin{block}{Correct Answer: B --- RPO}
		Recovery Point Objective is the maximum acceptable data loss measured in time. Four hours of data loss = RPO of 4 hours.
	\end{block}
\end{frame}

\begin{frame}{D2-Q449: RTO vs MTD}
	\begin{block}{Question}
		An application tolerates 4 hours of downtime but must be recovered within 6 hours. Which metrics are correct?
	\end{block}
	\begin{enumerate}[(A)]
		\item RPO 4h, RTO 6h
		\item MTD 4h, RTO 6h
		\item MTD 6h, RPO 4h
		\item RTO 4h, MTD 6h
	\end{enumerate}
\end{frame}
\begin{frame}{D2-Q449: Answer}
	\begin{block}{Correct Answer: D}
		RTO is the target recovery time (4h). MTD is the maximum tolerable downtime before severe impact (6h). RTO must be less than or equal to MTD.
	\end{block}
\end{frame}

\begin{frame}{D2-Q450: Risk Scenario Actors}
	\begin{block}{Question}
		Which risk scenario component has potential to generate internal or external threat?
	\end{block}
	\begin{enumerate}[(A)]
		\item Timing dimension
		\item Events
		\item Assets
		\item Actors
	\end{enumerate}
\end{frame}
\begin{frame}{D2-Q450: Answer}
	\begin{block}{Correct Answer: D}
		Actors—internal or external, human or non-human—have the potential to generate threats.
	\end{block}
\end{frame}

\begin{frame}{D2-Q451: Risk Scenario Requirements}
	\begin{block}{Question}
		What are the PRIMARY requirements for creating risk scenarios? (Choose two)
	\end{block}
	\begin{enumerate}[(A)]
		\item Potential threats and vulnerabilities
		\item Determination of value of asset/business process
		\item Determination of actors
		\item Determination of threat type
	\end{enumerate}
\end{frame}
\begin{frame}{D2-Q451: Answer}
	\begin{block}{Correct Answer: A and B}
		Risk scenarios require understanding of (1) threats and vulnerabilities that could cause loss and (2) the value of the asset or business process at risk.
	\end{block}
\end{frame}

\begin{frame}{D2-Q452: Effective Risk Assessment}
	\begin{block}{Question}
		Which is MOST effective in assessing business risk?
	\end{block}
	\begin{enumerate}[(A)]
		\item Use case analysis
		\item Business case analysis
		\item Risk scenarios
		\item Risk plan
	\end{enumerate}
\end{frame}
\begin{frame}{D2-Q452: Answer}
	\begin{block}{Correct Answer: C}
		Risk scenarios are the most effective technique for assessing business risk—they help determine likelihood and impact of identified risks.
	\end{block}
\end{frame}

\begin{frame}{D2-Q453: Risk Register Preparation}
	\begin{block}{Question}
		The preparation of a risk register begins in which risk management process?
	\end{block}
	\begin{enumerate}[(A)]
		\item Risk response planning
		\item Risk monitoring and control
		\item Risk management planning
		\item Risk identification
	\end{enumerate}
\end{frame}
\begin{frame}{D2-Q453: Answer}
	\begin{block}{Correct Answer: D}
		The risk register is initiated during risk identification, with description, cause, probability, impact, responses, owners, and status captured.
	\end{block}
\end{frame}

\begin{frame}{D2-Q454: BIA Primary Use}
	\begin{block}{Question}
		A Business Impact Analysis is PRIMARILY used to:
	\end{block}
	\begin{enumerate}[(A)]
		\item Estimate resources to resume operations
		\item Evaluate impact of disruption over time
		\item Calculate likelihood and impact of threats
		\item Evaluate high-level business requirements
	\end{enumerate}
\end{frame}
\begin{frame}{D2-Q454: Answer}
	\begin{block}{Correct Answer: B}
		A BIA evaluates impact of disruption over time on the enterprise's ability to operate and determines urgency of each business activity.
	\end{block}
\end{frame}

\begin{frame}{D2-Q455: Risk Scenarios Enable Assessment}
	\begin{block}{Question}
		Risk scenarios enable risk assessment because they:
	\end{block}
	\begin{enumerate}[(A)]
		\item Cover all potential risks
		\item Minimize need for quantitative methods
		\item Segregate IT from business risk
		\item Help estimate frequency and impact of risk
	\end{enumerate}
\end{frame}
\begin{frame}{D2-Q455: Answer}
	\begin{block}{Correct Answer: D}
		Risk scenarios help estimate frequency and impact—the two key dimensions of risk assessment.
	\end{block}
\end{frame}

\begin{frame}{D2-Q456: Risk Register for Scenarios}
	\begin{block}{Question}
		Which information in the risk register BEST helps develop proper risk scenarios? A list of:
	\end{block}
	\begin{enumerate}[(A)]
		\item Potential threats to assets
		\item Residual risk
		\item Accepted risk
		\item Security incidents
	\end{enumerate}
\end{frame}
\begin{frame}{D2-Q456: Answer}
	\begin{block}{Correct Answer: A}
		Potential threats to assets help isolate vulnerabilities and associated risks, which form the basis of risk scenarios.
	\end{block}
\end{frame}

\begin{frame}{D2-Q457: Identifying Control Deficiencies}
	\begin{block}{Question}
		Which BEST helps identify information systems control deficiencies?
	\end{block}
	\begin{enumerate}[(A)]
		\item Gap analysis
		\item Current IT risk profile
		\item IT controls framework
		\item Countermeasure analysis
	\end{enumerate}
\end{frame}
\begin{frame}{D2-Q457: Answer}
	\begin{block}{Correct Answer: A}
		Gap analysis compares desired control objectives with actual design and effectiveness to identify control deficiencies.
	\end{block}
\end{frame}

\begin{frame}{D2-Q458: Continuous Monitoring Performance}
	\begin{block}{Question}
		When assessing a critical application server, the MOST reliable assessment results come from:
	\end{block}
	\begin{enumerate}[(A)]
		\item Database auditing activation
		\item Documented performance objectives
		\item Continuous monitoring
		\item Security module documentation
	\end{enumerate}
\end{frame}
\begin{frame}{D2-Q458: Answer}
	\begin{block}{Correct Answer: C}
		Continuous monitoring tracks key performance metrics in real time and addresses issues with minimum impact.
	\end{block}
\end{frame}

\begin{frame}{D2-Q459: Risk Register Update Frequency}
	\begin{block}{Question}
		Which presents the GREATEST risk when updating the risk register? Updates are:
	\end{block}
	\begin{enumerate}[(A)]
		\item Carried out jointly
		\item Carried out after incidents
		\item Carried out annually
		\item Subject to CISO approval
	\end{enumerate}
\end{frame}
\begin{frame}{D2-Q459: Answer}
	\begin{block}{Correct Answer: C}
		Annual-only updates mean the register does not reflect current IT risk status. It should be updated continuously as risks change.
	\end{block}
\end{frame}

\begin{frame}{D2-Q460: IT Risk Measurement}
	\begin{block}{Question}
		IT risk is measured by its:
	\end{block}
	\begin{enumerate}[(A)]
		\item Damage to IT systems
		\item Impact on business operations
		\item Countermeasure cost
		\item Annual loss expectancy
	\end{enumerate}
\end{frame}
\begin{frame}{D2-Q460: Answer}
	\begin{block}{Correct Answer: B}
		IT risk is the business risk associated with IT; it is primarily measured by its impact on business operations.
	\end{block}
\end{frame}

\begin{frame}{D2-Q461: Quantitative Risk Analysis Description}
	\begin{block}{Question}
		Deriving likelihood and impact through statistical methods is BEST described as:
	\end{block}
	\begin{enumerate}[(A)]
		\item Quantitative risk analysis
		\item Risk scenario analysis
		\item Qualitative risk analysis
		\item Probabilistic risk assessment
	\end{enumerate}
\end{frame}
\begin{frame}{D2-Q461: Answer}
	\begin{block}{Correct Answer: A}
		Quantitative risk analysis derives probability and impact from statistical methods and data.
	\end{block}
\end{frame}

\begin{frame}{D2-Q462: Local Mitigation Responsibility}
	\begin{block}{Question}
		Local management has mitigated high-level risk. Who remains responsible?
	\end{block}
	\begin{enumerate}[(A)]
		\item Local management
		\item Corporate CRO
		\item Local purchasing manager
		\item Corporate management
	\end{enumerate}
\end{frame}
\begin{frame}{D2-Q462: Answer}
	\begin{block}{Correct Answer: D}
		Corporate management remains responsible for the risk even when responses are executed at a lower organizational level.
	\end{block}
\end{frame}

\begin{frame}{D2-Q463: Scenario Analysis for Likelihood}
	\begin{block}{Question}
		Which BEST estimates the likelihood of significant events affecting an enterprise?
	\end{block}
	\begin{enumerate}[(A)]
		\item Threat analysis
		\item Cost-benefit analysis
		\item Scenario analysis
		\item Countermeasure analysis
	\end{enumerate}
\end{frame}
\begin{frame}{D2-Q463: Answer}
	\begin{block}{Correct Answer: C}
		Scenario analysis with vulnerability analysis determines whether a risk is relevant and estimates likelihood of significant events.
	\end{block}
\end{frame}

\begin{frame}{D2-Q464: Improving Risk Decisions}
	\begin{block}{Question}
		Which BEST improves decision-making related to risk?
	\end{block}
	\begin{enumerate}[(A)]
		\item Documented risk register
		\item Risk awareness training
		\item Updated security policies
		\item Department-level accountability
	\end{enumerate}
\end{frame}
\begin{frame}{D2-Q464: Answer}
	\begin{block}{Correct Answer: A}
		A documented risk register captures relevant scenarios and prioritizes risk responses, improving decision quality.
	\end{block}
\end{frame}

\begin{frame}{D2-Q465: Independent Risk Review Purpose}
	\begin{block}{Question}
		The PRIMARY reason to have the risk management process reviewed by an independent professional is to:
	\end{block}
	\begin{enumerate}[(A)]
		\item Validate cost-effective solutions
		\item Validate internal control weaknesses
		\item Assess end-to-end process validity
		\item Assess whether risk profile and factors are defined properly
	\end{enumerate}
\end{frame}
\begin{frame}{D2-Q465: Answer}
	\begin{block}{Correct Answer: C}
		An independent reviewer can assess the entire risk management process end to end without bias.
	\end{block}
\end{frame}

% ============================================================
% DOMAIN 3 --- MORE RISK RESPONSE AND REPORTING (Q466–Q495)
% ============================================================
\section{Domain 3: More Risk Response and Reporting}

\begin{frame}{D3-Q466: Outsourced Noncore Greatest Concern}
	\begin{block}{Question}
		Which outcome of outsourcing noncore processes is of GREATEST concern?
	\end{block}
	\begin{enumerate}[(A)]
		\item Total cost exceeds projections
		\item Internal IS experience is lost
		\item Vendor employees are disloyal
		\item Sensitive data is subcontracted by the vendor
	\end{enumerate}
\end{frame}
\begin{frame}{D3-Q466: Answer}
	\begin{block}{Correct Answer: D}
		Subcontracting sensitive data increases cession of control—the greatest concern in third-party relationships.
	\end{block}
\end{frame}

\begin{frame}{D3-Q467: Business Case for Risk Treatment}
	\begin{block}{Question}
		When proposing a risk treatment, a risk practitioner PRIMARILY uses a:
	\end{block}
	\begin{enumerate}[(A)]
		\item Technical evaluation report
		\item Business case
		\item Vulnerability report
		\item Budgetary requirements
	\end{enumerate}
\end{frame}
\begin{frame}{D3-Q467: Answer}
	\begin{block}{Correct Answer: B}
		A business case documents the risk assessment, control costs, and expected benefits—essential for approval.
	\end{block}
\end{frame}

\begin{frame}{D3-Q468: Most Important for Effective Risk Management}
	\begin{block}{Question}
		Which is MOST important for effective risk management?
	\end{block}
	\begin{enumerate}[(A)]
		\item Assigning risk owners
		\item Ensuring regulatory compliance
		\item Integrating risk into operations
		\item Implementing risk avoidance
	\end{enumerate}
\end{frame}
\begin{frame}{D3-Q468: Answer}
	\begin{block}{Correct Answer: A}
		Assigning risk to individual owners maximizes accountability—the foundation of effective risk management.
	\end{block}
\end{frame}

\begin{frame}{D3-Q469: Most Important Outsourcing Question}
	\begin{block}{Question}
		What is the MOST important question about outsourcing arrangements?
	\end{block}
	\begin{enumerate}[(A)]
		\item Are security exception procedures in place?
		\item Are SLAs being met?
		\item Is the provider ISO certified?
		\item Are specific security controls mandated in the contract?
	\end{enumerate}
\end{frame}
\begin{frame}{D3-Q469: Answer}
	\begin{block}{Correct Answer: D}
		Without explicit security requirements in the contract, the enterprise has no assurance that the provider will comply with specific controls.
	\end{block}
\end{frame}

\begin{frame}{D3-Q470: Basis for Risk Ownership}
	\begin{block}{Question}
		Which activity provides the BEST basis for establishing risk ownership?
	\end{block}
	\begin{enumerate}[(A)]
		\item Documenting interdependencies
		\item Mapping risk to specific business processes
		\item Reviewing RACI charts
		\item Distributing risk among business owners
	\end{enumerate}
\end{frame}
\begin{frame}{D3-Q470: Answer}
	\begin{block}{Correct Answer: B}
		Mapping risk to a specific business process identifies the process owner—the natural risk owner.
	\end{block}
\end{frame}

\begin{frame}{D3-Q471: Risk Response Report}
	\begin{block}{Question}
		A risk response report includes recommendations for:
	\end{block}
	\begin{enumerate}[(A)]
		\item Acceptance
		\item Assessment
		\item Evaluation
		\item Quantification
	\end{enumerate}
\end{frame}
\begin{frame}{D3-Q471: Answer}
	\begin{block}{Correct Answer: A}
		Risk response reports recommend among the four response options: accept, mitigate, transfer, avoid.
	\end{block}
\end{frame}

\begin{frame}{D3-Q472: Reducing to Acceptable Risk}
	\begin{block}{Question}
		Which is MOST likely reduced to achieve acceptable risk?
	\end{block}
	\begin{enumerate}[(A)]
		\item Risk appetite
		\item Control risk
		\item Residual risk
		\item Inherent risk
	\end{enumerate}
\end{frame}
\begin{frame}{D3-Q472: Answer}
	\begin{block}{Correct Answer: C}
		Residual risk—the risk remaining after controls—is what must be reduced to within appetite.
	\end{block}
\end{frame}

\begin{frame}{D3-Q473: Denial-of-Service Decision}
	\begin{block}{Question}
		A financial institution decides not to address a DoS vulnerability. The MOST likely reason is:
	\end{block}
	\begin{enumerate}[(A)]
		\item Countermeasure is too complex
		\item Sufficient safeguards exist
		\item Likelihood is unknown
		\item Cost of countermeasure outweighs potential loss
	\end{enumerate}
\end{frame}
\begin{frame}{D3-Q473: Answer}
	\begin{block}{Correct Answer: D}
		The enterprise accepted the risk because the cost of mitigation exceeded the potential loss.
	\end{block}
\end{frame}

\begin{frame}{D3-Q474: Two-Factor Cost-Benefit}
	\begin{block}{Question}
		Which is MOST relevant to a two-factor authentication cost-benefit analysis?
	\end{block}
	\begin{enumerate}[(A)]
		\item Approved project budget
		\item Incident frequency
		\item Annual loss expectancy
		\item Total cost of ownership
	\end{enumerate}
\end{frame}
\begin{frame}{D3-Q474: Answer}
	\begin{block}{Correct Answer: D}
		Total cost of ownership establishes the full life-cycle cost baseline required for cost-benefit comparison.
	\end{block}
\end{frame}

\begin{frame}{D3-Q475: Outsourcing Contract Most Important}
	\begin{block}{Question}
		What is the MOST important part of an outsourcing contract?
	\end{block}
	\begin{enumerate}[(A)]
		\item Right to audit
		\item Provisions to assess compliance
		\item Incident notification procedures
		\item Encryption requirements
	\end{enumerate}
\end{frame}
\begin{frame}{D3-Q475: Answer}
	\begin{block}{Correct Answer: B}
		Without compliance assessment provisions, the enterprise cannot verify that the supplier meets its security obligations.
	\end{block}
\end{frame}

\begin{frame}{D3-Q476: Ambiguous SLA Risk}
	\begin{block}{Question}
		What is the GREATEST risk when engaging a cloud provider?
	\end{block}
	\begin{enumerate}[(A)]
		\item Same geographical location
		\item Ambiguous SLA
		\item Missing references
		\item Reliance on third-party audits
	\end{enumerate}
\end{frame}
\begin{frame}{D3-Q476: Answer}
	\begin{block}{Correct Answer: B}
		Ambiguous SLAs make it difficult to determine whether the provider meets its obligations—leaving the enterprise without recourse.
	\end{block}
\end{frame}

\begin{frame}{D3-Q477: Insurance for Unmitigable Risk}
	\begin{block}{Question}
		When a risk cannot be sufficiently mitigated, the BEST option to protect financial impact is:
	\end{block}
	\begin{enumerate}[(A)]
		\item Insurance
		\item Updating the register
		\item Staff training
		\item Outsourcing
	\end{enumerate}
\end{frame}
\begin{frame}{D3-Q477: Answer}
	\begin{block}{Correct Answer: A}
		Insurance transfers financial risk to the insurer, providing monetary compensation for impact.
	\end{block}
\end{frame}

\begin{frame}{D3-Q478: Cost Exceeds Benefit}
	\begin{block}{Question}
		When mitigation cost exceeds benefit, the BEST risk response is:
	\end{block}
	\begin{enumerate}[(A)]
		\item Treat
		\item Terminate
		\item Accept
		\item Transfer
	\end{enumerate}
\end{frame}
\begin{frame}{D3-Q478: Answer}
	\begin{block}{Correct Answer: C}
		When control cost exceeds expected loss, accepting the risk is economically rational.
	\end{block}
\end{frame}

\begin{frame}{D3-Q479: Risk Scenarios for Continuity}
	\begin{block}{Question}
		Which activity MOST benefits a continuity improvement initiative?
	\end{block}
	\begin{enumerate}[(A)]
		\item Develop KRIs
		\item Determine control ownership
		\item Update the risk register
		\item Develop IT risk scenarios
	\end{enumerate}
\end{frame}
\begin{frame}{D3-Q479: Answer}
	\begin{block}{Correct Answer: D}
		Risk scenarios link IT failures to business impact, informing continuity planning and prioritization.
	\end{block}
\end{frame}

\begin{frame}{D3-Q480: Reciprocal Agreement}
	\begin{block}{Question}
		A reciprocal agreement with another company for equipment is an example of risk:
	\end{block}
	\begin{enumerate}[(A)]
		\item Transfer
		\item Avoidance
		\item Acceptance
		\item Mitigation
	\end{enumerate}
\end{frame}
\begin{frame}{D3-Q480: Answer}
	\begin{block}{Correct Answer: D}
		Reciprocal agreements reduce impact when an event occurs—a form of risk mitigation.
	\end{block}
\end{frame}

\begin{frame}{D3-Q481: Selecting Risk Response}
	\begin{block}{Question}
		Which would BEST help an enterprise select an appropriate risk response?
	\end{block}
	\begin{enumerate}[(A)]
		\item Change analysis
		\item Transferable risk analysis
		\item Likelihood and impact analysis
		\item Control cost and benefit analysis
	\end{enumerate}
\end{frame}
\begin{frame}{D3-Q481: Answer}
	\begin{block}{Correct Answer: D}
		Cost-benefit analysis determines whether mitigation can bring risk to acceptable levels economically.
	\end{block}
\end{frame}

\begin{frame}{D3-Q482: Optimal Security Investment}
	\begin{block}{Question}
		Which leads to optimal return on security investment?
	\end{block}
	\begin{enumerate}[(A)]
		\item Maximum protection on all assets
		\item Focus on important assets and protect appropriately
		\item Minimum protection on all assets
		\item Invest only after major incidents
	\end{enumerate}
\end{frame}
\begin{frame}{D3-Q482: Answer}
	\begin{block}{Correct Answer: B}
		Optimizing ROI means protecting critical assets appropriately while avoiding overprotection of less critical ones.
	\end{block}
\end{frame}

\begin{frame}{D3-Q483: Compromised Admin Account}
	\begin{block}{Question}
		Which control reduces impact of a compromised administrator account?
	\end{block}
	\begin{enumerate}[(A)]
		\item Multifactor authentication
		\item Audit logging
		\item Least privilege
		\item Password policy
	\end{enumerate}
\end{frame}
\begin{frame}{D3-Q483: Answer}
	\begin{block}{Correct Answer: C}
		Least privilege limits the scope of damage even when credentials are compromised.
	\end{block}
\end{frame}

\begin{frame}{D3-Q484: Critical Element of Risk Program}
	\begin{block}{Question}
		Which is critical for an effective risk management program?
	\end{block}
	\begin{enumerate}[(A)]
		\item Response strategy
		\item Register content
		\item Profile content
		\item Risk owners and accountability
	\end{enumerate}
\end{frame}
\begin{frame}{D3-Q484: Answer}
	\begin{block}{Correct Answer: D}
		Risk owners are essential—they make informed, cost-effective decisions about controls for their owned risks.
	\end{block}
\end{frame}

\begin{frame}{D3-Q485: Anti-Malware Cost Exceeds Loss}
	\begin{block}{Question}
		When anti-malware cost exceeds loss expectancy, the MOST viable response is:
	\end{block}
	\begin{enumerate}[(A)]
		\item Eliminate risk
		\item Accept risk
		\item Transfer risk
		\item Mitigate risk
	\end{enumerate}
\end{frame}
\begin{frame}{D3-Q485: Answer}
	\begin{block}{Correct Answer: B}
		When response cost exceeds expected loss, risk acceptance is most economically viable.
	\end{block}
\end{frame}

\begin{frame}{D3-Q486: Risk Avoidance Behavior}
	\begin{block}{Question}
		Which is the BEST example of risk avoidance?
	\end{block}
	\begin{enumerate}[(A)]
		\item Taking no action
		\item Outsourcing the process
		\item Insuring against an event
		\item Exiting the activity that creates risk
	\end{enumerate}
\end{frame}
\begin{frame}{D3-Q486: Answer}
	\begin{block}{Correct Answer: D}
		Avoidance eliminates risk by exiting the activity—not by retaining, transferring, or mitigating it.
	\end{block}
\end{frame}

\begin{frame}{D3-Q487: Critical Dependencies in Treatment}
	\begin{block}{Question}
		Which ensures critical dependencies are addressed in the risk treatment plan?
	\end{block}
	\begin{enumerate}[(A)]
		\item Treat all possible risk
		\item Verify business objectives via top-down process review
		\item Treat each risk independently
		\item Verify objectives via bottom-up review
	\end{enumerate}
\end{frame}
\begin{frame}{D3-Q487: Answer}
	\begin{block}{Correct Answer: B}
		A top-down approach driven by business objectives ensures critical dependencies are addressed holistically.
	\end{block}
\end{frame}

\begin{frame}{D3-Q488: Controls to Mitigate Risk}
	\begin{block}{Question}
		A CISO recommends anti-malware and similar controls. Which approach is being applied?
	\end{block}
	\begin{enumerate}[(A)]
		\item Transference
		\item Mitigation
		\item Acceptance
		\item Avoidance
	\end{enumerate}
\end{frame}
\begin{frame}{D3-Q488: Answer}
	\begin{block}{Correct Answer: B}
		Implementing controls to reduce risk to acceptable levels is risk mitigation.
	\end{block}
\end{frame}

\begin{frame}{D3-Q489: Best Outsourcing Mitigation}
	\begin{block}{Question}
		Which practice BEST mitigates outsourcing risk?
	\end{block}
	\begin{enumerate}[(A)]
		\item Audits to verify compliance
		\item Vendor awareness training
		\item Retaining copies of data
		\item Reviewing vendor financials
	\end{enumerate}
\end{frame}
\begin{frame}{D3-Q489: Answer}
	\begin{block}{Correct Answer: A}
		Audits verify contract compliance and are the strongest mitigation for outsourcing risk.
	\end{block}
\end{frame}

\begin{frame}{D3-Q490: Best Risk Transfer}
	\begin{block}{Question}
		Which risk is BEST addressed through transfer?
	\end{block}
	\begin{enumerate}[(A)]
		\item Old fire suppression system
		\item Disgruntled employee sabotage
		\item Disgruntled employee theft
		\item Building in a 100-year flood plain
	\end{enumerate}
\end{frame}
\begin{frame}{D3-Q490: Answer}
	\begin{block}{Correct Answer: D}
		Flood insurance transfers the financial risk of a flood—a classic risk transfer scenario.
	\end{block}
\end{frame}

\begin{frame}{D3-Q491: Finalizing Risk Treatment Plan}
	\begin{block}{Question}
		Which would BEST help finalize the risk treatment plan?
	\end{block}
	\begin{enumerate}[(A)]
		\item Vulnerability analysis
		\item Impact analysis
		\item Cost-benefit analysis
		\item SWOT analysis
	\end{enumerate}
\end{frame}
\begin{frame}{D3-Q491: Answer}
	\begin{block}{Correct Answer: C}
		Cost-benefit analysis confirms that control benefits exceed costs, finalizing the treatment decision.
	\end{block}
\end{frame}

\begin{frame}{D3-Q492: Greatest Treatment Plan Benefit}
	\begin{block}{Question}
		The GREATEST benefit of a risk treatment plan is to:
	\end{block}
	\begin{enumerate}[(A)]
		\item Reduce impact and likelihood of occurrence
		\item Identify transferable risk
		\item Exploit risk for preparedness
		\item Enhance risk appetite
	\end{enumerate}
\end{frame}
\begin{frame}{D3-Q492: Answer}
	\begin{block}{Correct Answer: A}
		Treatment reduces the impact and likelihood of a risk event—its core purpose.
	\end{block}
\end{frame}

\begin{frame}{D3-Q493: Validating Patching Program}
	\begin{block}{Question}
		What is the BEST method to validate the patching program?
	\end{block}
	\begin{enumerate}[(A)]
		\item Penetration testing
		\item Risk identification
		\item Vulnerability scanning
		\item Request-for-change review
	\end{enumerate}
\end{frame}
\begin{frame}{D3-Q493: Answer}
	\begin{block}{Correct Answer: C}
		Vulnerability scans reveal whether patches have been installed on a timely basis.
	\end{block}
\end{frame}

\begin{frame}{D3-Q494: Most Important IS Control Design Factor}
	\begin{block}{Question}
		What is the MOST important factor when designing IS controls in a complex environment?
	\end{block}
	\begin{enumerate}[(A)]
		\item Development methodology
		\item Scalability
		\item Platform interfaces
		\item Stakeholder requirements
	\end{enumerate}
\end{frame}
\begin{frame}{D3-Q494: Answer}
	\begin{block}{Correct Answer: D}
		Controls must address stakeholder and business requirements to advance enterprise interests.
	\end{block}
\end{frame}

\begin{frame}{D3-Q495: EFT Testing Priority}
	\begin{block}{Question}
		After major EFT enhancements, the risk practitioner would MOST benefit from a test that:
	\end{block}
	\begin{enumerate}[(A)]
		\item Identifies new gaps in security
		\item Verifies recovery capability
		\item Ensures performance to expectation
		\item Ensures volume support
	\end{enumerate}
\end{frame}
\begin{frame}{D3-Q495: Answer}
	\begin{block}{Correct Answer: A}
		For a financial transfer system, security is the highest priority for risk practitioners; testing should reveal new security gaps.
	\end{block}
\end{frame}

% ============================================================
% DOMAIN 4 --- MORE IT AND SECURITY (Q496–Q515)
% ============================================================
\section{Domain 4: More IT and Security}

\begin{frame}{D4-Q496: Log Management Concern}
	\begin{block}{Question}
		Which is the PRIMARY purpose of log aggregation?
	\end{block}
	\begin{enumerate}[(A)]
		\item Identify and analyze important events
		\item Synchronize clocks
		\item Eliminate logging
		\item Centralize storage only
	\end{enumerate}
\end{frame}
\begin{frame}{D4-Q496: Answer}
	\begin{block}{Correct Answer: A}
		Log aggregation helps identify and analyze important events and potential incidents across systems.
	\end{block}
\end{frame}

\begin{frame}{D4-Q497: SIEM Function}
	\begin{block}{Question}
		SIEM systems primarily:
	\end{block}
	\begin{enumerate}[(A)]
		\item Correlate data from multiple sources
		\item Store backup data
		\item Manage user identities
		\item Encrypt network traffic
	\end{enumerate}
\end{frame}
\begin{frame}{D4-Q497: Answer}
	\begin{block}{Correct Answer: A}
		SIEM tools capture and correlate events from multiple sources to detect security incidents.
	\end{block}
\end{frame}

\begin{frame}{D4-Q498: Log Review Challenge}
	\begin{block}{Question}
		A common challenge with log review is:
	\end{block}
	\begin{enumerate}[(A)]
		\item Too much data to analyze effectively
		\item Not enough data collected
		\item Logs too expensive
		\item Logs not required
	\end{enumerate}
\end{frame}
\begin{frame}{D4-Q498: Answer}
	\begin{block}{Correct Answer: A}
		Data volume, difficulty searching, misconfiguration, and pre-read modification are common log review challenges.
	\end{block}
\end{frame}

\begin{frame}{D4-Q499: Control Monitoring Purpose}
	\begin{block}{Question}
		The purpose of control monitoring is to verify:
	\end{block}
	\begin{enumerate}[(A)]
		\item Whether the control is working
		\item Whether the control is effectively addressing risk
		\item Whether staff follows procedures
		\item Whether logs are reviewed
	\end{enumerate}
\end{frame}
\begin{frame}{D4-Q499: Answer}
	\begin{block}{Correct Answer: B}
		Control monitoring verifies the control is effectively addressing risk—not merely that it operates.
	\end{block}
\end{frame}

\begin{frame}{D4-Q500: KRI Sensitivity}
	\begin{block}{Question}
		A KRI that generates too many false positives should be adjusted for:
	\end{block}
	\begin{enumerate}[(A)]
		\item Sensitivity
		\item Timeliness
		\item Frequency
		\item Reliability
	\end{enumerate}
\end{frame}
\begin{frame}{D4-Q500: Answer}
	\begin{block}{Correct Answer: A}
		Sensitivity thresholds control how often a KRI alerts; too many false positives indicates the threshold is too sensitive.
	\end{block}
\end{frame}

\begin{frame}{D4-Q501: KRI vs KPI}
	\begin{block}{Question}
		Which statement BEST distinguishes KRIs from KPIs?
	\end{block}
	\begin{enumerate}[(A)]
		\item They measure the same thing
		\item KPIs identify underperformance; KRIs provide early warnings
		\item KRIs replace KPIs in mature organizations
		\item They are unrelated
	\end{enumerate}
\end{frame}
\begin{frame}{D4-Q501: Answer}
	\begin{block}{Correct Answer: B}
		KPIs measure performance against objectives; KRIs provide early warnings of increasing risk. Both are complementary.
	\end{block}
\end{frame}

\begin{frame}{D4-Q502: KRI Evaluation Frequency}
	\begin{block}{Question}
		KRIs should be evaluated:
	\end{block}
	\begin{enumerate}[(A)]
		\item Annually
		\item Only when incidents occur
		\item On a regular basis
		\item Never after implementation
	\end{enumerate}
\end{frame}
\begin{frame}{D4-Q502: Answer}
	\begin{block}{Correct Answer: C}
		Because the internal and external environments change continuously, KRIs must be evaluated regularly to remain aligned with appetite.
	\end{block}
\end{frame}

\begin{frame}{D4-Q503: Control Effectiveness Verification}
	\begin{block}{Question}
		One way to verify control effectiveness is by determining:
	\end{block}
	\begin{enumerate}[(A)]
		\item Its reliability
		\item Whether it is preventive or detective
		\item Capability of failure notification
		\item Test results of intended objectives
	\end{enumerate}
\end{frame}
\begin{frame}{D4-Q503: Answer}
	\begin{block}{Correct Answer: D}
		Effectiveness is verified by testing whether the control achieves its intended objectives.
	\end{block}
\end{frame}

\begin{frame}{D4-Q504: KRI Reliability}
	\begin{block}{Question}
		Which BEST describes KRI reliability?
	\end{block}
	\begin{enumerate}[(A)]
		\item It operates within thresholds
		\item It tests the target at intervals
		\item It flags exceptions every time they occur
		\item It initiates corrective action
	\end{enumerate}
\end{frame}
\begin{frame}{D4-Q504: Answer}
	\begin{block}{Correct Answer: C}
		A reliable KRI consistently flags exceptions when they occur, indicating a strong correlation with the risk.
	\end{block}
\end{frame}

\begin{frame}{D4-Q505: SIEM for Incident Source}
	\begin{block}{Question}
		Which system BEST provides information for investigating the root cause of a security incident?
	\end{block}
	\begin{enumerate}[(A)]
		\item IDS
		\item IPS
		\item Configuration management
		\item Event management (SIEM)
	\end{enumerate}
\end{frame}
\begin{frame}{D4-Q505: Answer}
	\begin{block}{Correct Answer: D}
		SIEM merges logs from multiple systems, enabling full data set analysis for root cause investigation.
	\end{block}
\end{frame}

\begin{frame}{D4-Q506: Best Control for Unauthorized Access}
	\begin{block}{Question}
		The BEST control to prevent unauthorized access is user:
	\end{block}
	\begin{enumerate}[(A)]
		\item Accountability
		\item Authentication
		\item Identification
		\item Access rules
	\end{enumerate}
\end{frame}
\begin{frame}{D4-Q506: Answer}
	\begin{block}{Correct Answer: B}
		Authentication verifies user identity before granting access—preventing unauthorized users.
	\end{block}
\end{frame}

\begin{frame}{D4-Q507: Need-to-Know Access}
	\begin{block}{Question}
		The BEST control for protecting sensitive information from unauthorized access is:
	\end{block}
	\begin{enumerate}[(A)]
		\item Challenge-response system
		\item Periodic password changes
		\item Monitoring failed logins
		\item Need-to-know access
	\end{enumerate}
\end{frame}
\begin{frame}{D4-Q507: Answer}
	\begin{block}{Correct Answer: D}
		Need-to-know access embodies least privilege and segregation of duties—fundamental controls.
	\end{block}
\end{frame}

\begin{frame}{D4-Q508: Best Defense Against Phishing}
	\begin{block}{Question}
		The BEST defense against phishing attacks is:
	\end{block}
	\begin{enumerate}[(A)]
		\item IDS
		\item Spam filters
		\item End-user awareness
		\item Application hardening
	\end{enumerate}
\end{frame}
\begin{frame}{D4-Q508: Answer}
	\begin{block}{Correct Answer: C}
		Phishing is a social engineering attack; user awareness training is the best defense.
	\end{block}
\end{frame}

\begin{frame}{D4-Q509: Hashing Purpose}
	\begin{block}{Question}
		The primary purpose of hashing is to:
	\end{block}
	\begin{enumerate}[(A)]
		\item Provide confidentiality
		\item Verify data integrity
		\item Provide authorization
		\item Facilitate key exchange
	\end{enumerate}
\end{frame}
\begin{frame}{D4-Q509: Answer}
	\begin{block}{Correct Answer: B}
		Hashing generates a fixed-length value used to verify that data has not been altered.
	\end{block}
\end{frame}

\begin{frame}{D4-Q510: Digital Signature Provides}
	\begin{block}{Question}
		Which does NOT come from a digital signature?
	\end{block}
	\begin{enumerate}[(A)]
		\item Integrity
		\item Nonrepudiation
		\item Authenticity
		\item Confidentiality
	\end{enumerate}
\end{frame}
\begin{frame}{D4-Q510: Answer}
	\begin{block}{Correct Answer: D}
		Digital signatures provide integrity, nonrepudiation, and authenticity—not confidentiality.
	\end{block}
\end{frame}

\begin{frame}{D4-Q511: Encrypt for Receiver Only}
	\begin{block}{Question}
		To ensure only the receiver can decrypt, the sender uses:
	\end{block}
	\begin{enumerate}[(A)]
		\item Receiver's public key
		\item Sender's public key
		\item Receiver's private key
		\item Sender's private key
	\end{enumerate}
\end{frame}
\begin{frame}{D4-Q511: Answer}
	\begin{block}{Correct Answer: A}
		Encrypting with the receiver's public key ensures only the receiver can decrypt with their private key.
	\end{block}
\end{frame}

\begin{frame}{D4-Q512: Signed and Encrypted Provides}
	\begin{block}{Question}
		Which cryptographic mechanism provides confidentiality, integrity, and nonrepudiation?
	\end{block}
	\begin{enumerate}[(A)]
		\item Signed only
		\item Encrypted only
		\item Signed and encrypted
		\item Signed and hashed
	\end{enumerate}
\end{frame}
\begin{frame}{D4-Q512: Answer}
	\begin{block}{Correct Answer: C}
		Encryption provides confidentiality; digital signature provides integrity and nonrepudiation. Together they cover all three.
	\end{block}
\end{frame}

\begin{frame}{D4-Q513: PKI Single Point of Failure}
	\begin{block}{Question}
		Which is a single point of failure in a PKI ecosystem?
	\end{block}
	\begin{enumerate}[(A)]
		\item Certificate
		\item Certificate Authority (CA)
		\item Requester
		\item Browser
	\end{enumerate}
\end{frame}
\begin{frame}{D4-Q513: Answer}
	\begin{block}{Correct Answer: B}
		If the CA's private key is compromised, confidence in all certificates it issued collapses.
	\end{block}
\end{frame}

\begin{frame}{D4-Q514: Security Awareness Primary Goal}
	\begin{block}{Question}
		The most important goal of security awareness training is:
	\end{block}
	\begin{enumerate}[(A)]
		\item Train on latest threats
		\item Demonstrate audit compliance
		\item Create a security-conscious culture
		\item Upskill employees
	\end{enumerate}
\end{frame}
\begin{frame}{D4-Q514: Answer}
	\begin{block}{Correct Answer: C}
		Creating a security-conscious culture is the primary, lasting goal of awareness training.
	\end{block}
\end{frame}

\begin{frame}{D4-Q515: Data Security vs Privacy}
	\begin{block}{Question}
		Which statement about data security and data privacy is TRUE?
	\end{block}
	\begin{enumerate}[(A)]
		\item They are identical
		\item Security protects data; privacy governs its proper use
		\item Privacy is more important than security
		\item Security applies only to personal data
	\end{enumerate}
\end{frame}
\begin{frame}{D4-Q515: Answer}
	\begin{block}{Correct Answer: B}
		Security protects data from unauthorized access; privacy governs appropriate use of personal information. Strong security does not guarantee privacy.
	\end{block}
\end{frame}


\section{Domain 1: Additional Governance}

\begin{frame}{D1-Q316: Risk and Control Ownership}
	\begin{block}{Question}
		Which of the following is MOST important when assigning risk and control ownership in an enterprise?
	\end{block}
	\begin{enumerate}[(A)]
		\item The number of controls each owner manages
		\item Clear definition of roles and accountability
		\item The technical expertise of each owner
		\item The organizational level of each owner
	\end{enumerate}
\end{frame}
\begin{frame}{D1-Q316: Answer}
	\begin{block}{Correct Answer: B}
		Clear definition of roles and accountability is the most important factor. Without unambiguous ownership, controls may be duplicated, neglected, or contested, and accountability for risk cannot be enforced.
	\end{block}
\end{frame}

\begin{frame}{D1-Q317: Board Oversight of IT Risk}
	\begin{block}{Question}
		Which mechanism BEST enables a board of directors to oversee IT risk effectively?
	\end{block}
	\begin{enumerate}[(A)]
		\item Receiving detailed technical vulnerability reports
		\item Periodic reporting of aggregated risk against appetite
		\item Delegating all IT risk decisions to the CIO
		\item Reviewing individual firewall configurations
	\end{enumerate}
\end{frame}
\begin{frame}{D1-Q317: Answer}
	\begin{block}{Correct Answer: B}
		Boards need aggregated risk information aligned to risk appetite and tolerance, not technical detail. This allows them to provide direction and oversight without micromanaging operations.
	\end{block}
\end{frame}

\begin{frame}{D1-Q318: Risk Governance Principles}
	\begin{block}{Question}
		Which principle is MOST fundamental to effective risk governance?
	\end{block}
	\begin{enumerate}[(A)]
		\item Risk is managed primarily by the internal audit function
		\item Risk decisions are made at the lowest possible cost
		\item Risk oversight is independent from risk-taking activities
		\item All risk must be eliminated before business proceeds
	\end{enumerate}
\end{frame}
\begin{frame}{D1-Q318: Answer}
	\begin{block}{Correct Answer: C}
		Independence of oversight from risk-taking activities ensures objectivity. This is the foundation of the three lines of defense model and prevents conflicts of interest.
	\end{block}
\end{frame}

\begin{frame}{D1-Q319: Enterprise Risk Culture Indicators}
	\begin{block}{Question}
		Which is the BEST indicator of a healthy enterprise risk culture?
	\end{block}
	\begin{enumerate}[(A)]
		\item Few risks are reported because issues are resolved quietly
		\item Employees escalate concerns without fear of retaliation
		\item Risk assessments are performed only when required
		\item Senior management rarely discusses risk
	\end{enumerate}
\end{frame}
\begin{frame}{D1-Q319: Answer}
	\begin{block}{Correct Answer: B}
		A healthy risk culture encourages open escalation and psychological safety. Fear of retaliation is a hallmark of a blame culture that suppresses risk reporting.
	\end{block}
\end{frame}

\begin{frame}{D1-Q320: Risk Appetite Statement Content}
	\begin{block}{Question}
		Which element is MOST essential in a risk appetite statement?
	\end{block}
	\begin{enumerate}[(A)]
		\item A list of all identified risks
		\item The amount and type of risk the enterprise is willing to accept
		\item Detailed control procedures for each business unit
		\item The names of all risk owners
	\end{enumerate}
\end{frame}
\begin{frame}{D1-Q320: Answer}
	\begin{block}{Correct Answer: B}
		A risk appetite statement articulates the amount and type of risk the enterprise is willing to accept in pursuit of its objectives. It guides decision-making at all levels.
	\end{block}
\end{frame}

\begin{frame}{D1-Q321: Risk Culture and Incentives}
	\begin{block}{Question}
		How do incentive structures MOST influence enterprise risk culture?
	\end{block}
	\begin{enumerate}[(A)]
		\item They have no effect on how employees treat risk
		\item They align employee behavior with desired risk-taking boundaries
		\item They eliminate the need for formal risk policies
		\item They guarantee that no risk events will occur
	\end{enumerate}
\end{frame}
\begin{frame}{D1-Q321: Answer}
	\begin{block}{Correct Answer: B}
		Incentives that reward appropriate risk-taking and penalize reckless behavior align employee actions with the enterprise's risk appetite. Misaligned incentives can undermine even well-designed controls.
	\end{block}
\end{frame}

\begin{frame}{D1-Q322: Risk Governance Maturity}
	\begin{block}{Question}
		Which characteristic BEST indicates a mature risk governance structure?
	\end{block}
	\begin{enumerate}[(A)]
		\item Risk is managed reactively after incidents occur
		\item Risk considerations are integrated into strategic planning
		\item Risk management is performed only by the risk department
		\item Risk reports are produced annually for compliance
	\end{enumerate}
\end{frame}
\begin{frame}{D1-Q322: Answer}
	\begin{block}{Correct Answer: B}
		Mature governance integrates risk into strategy, budgeting, and operations. Reactive or siloed risk management indicates lower maturity.
	\end{block}
\end{frame}

\begin{frame}{D1-Q323: Communicating Risk Appetite}
	\begin{block}{Question}
		What is the BEST way to ensure risk appetite is understood throughout the enterprise?
	\end{block}
	\begin{enumerate}[(A)]
		\item Publish the appetite statement only on the intranet
		\item Translate appetite into specific tolerances and metrics for each business unit
		\item Discuss appetite only at board meetings
		\item Include appetite in the annual audit report
	\end{enumerate}
\end{frame}
\begin{frame}{D1-Q323: Answer}
	\begin{block}{Correct Answer: B}
		Translating high-level appetite into concrete tolerances, thresholds, and metrics makes it actionable for each business unit. A statement alone is insufficient.
	\end{block}
\end{frame}

\begin{frame}{D1-Q324: Risk Governance Failures}
	\begin{block}{Question}
		Which is the MOST common root cause of risk governance failures?
	\end{block}
	\begin{enumerate}[(A)]
		\item Insufficient technical tools
		\item Lack of clear accountability and oversight
		\item Too many risk assessments
		\item Excessive regulatory requirements
	\end{enumerate}
\end{frame}
\begin{frame}{D1-Q324: Answer}
	\begin{block}{Correct Answer: B}
		Governance failures most often stem from unclear accountability, inadequate oversight, and poorly defined roles—not from tooling or assessment frequency.
	\end{block}
\end{frame}

\begin{frame}{D1-Q325: Board Risk Committee}
	\begin{block}{Question}
		What is the PRIMARY role of a board-level risk committee?
	\end{block}
	\begin{enumerate}[(A)]
		\item Execute day-to-day risk management activities
		\item Oversee the enterprise risk management framework and risk appetite
		\item Conduct detailed technical risk assessments
		\item Approve individual risk treatment plans
	\end{enumerate}
\end{frame}
\begin{frame}{D1-Q325: Answer}
	\begin{block}{Correct Answer: B}
		A board risk committee provides oversight of the ERM framework, risk appetite, and key risk exposures. Execution and technical assessment are delegated to management.
	\end{block}
\end{frame}

\begin{frame}{D1-Q326: Risk Reporting to the Board}
	\begin{block}{Question}
		Which characteristic makes risk reporting MOST useful to the board?
	\end{block}
	\begin{enumerate}[(A)]
		\item Inclusion of every identified risk in detail
		\item Focus on aggregated exposures against appetite with clear trends
		\item Use of technical jargon to demonstrate expertise
		\item Frequency of at least weekly
	\end{enumerate}
\end{frame}
\begin{frame}{D1-Q326: Answer}
	\begin{block}{Correct Answer: B}
		Boards need aggregated, trend-oriented views against appetite—not exhaustive detail. This enables strategic oversight and informed decision-making.
	\end{block}
\end{frame}

\begin{frame}{D1-Q327: Risk Ownership Delegation}
	\begin{block}{Question}
		Which statement about risk ownership delegation is TRUE?
	\end{block}
	\begin{enumerate}[(A)]
		\item Accountability can be delegated along with responsibility
		\item Responsibility can be delegated, but accountability cannot
		\item Neither responsibility nor accountability can be delegated
		\item Both responsibility and accountability are optional
	\end{enumerate}
\end{frame}
\begin{frame}{D1-Q327: Answer}
	\begin{block}{Correct Answer: B}
		Responsibility for tasks can be delegated, but accountability for outcomes remains with the original owner. This is a core governance principle.
	\end{block}
\end{frame}

\begin{frame}{D1-Q328: Risk Governance in M\&A}
	\begin{block}{Question}
		What is the FIRST risk governance consideration during a merger or acquisition?
	\end{block}
	\begin{enumerate}[(A)]
		\item Harmonizing technical platforms
		\item Assessing the target's risk culture and governance maturity
		\item Reducing headcount
		\item Standardizing job titles
	\end{enumerate}
\end{frame}
\begin{frame}{D1-Q328: Answer}
	\begin{block}{Correct Answer: B}
		Understanding the target's risk culture and governance maturity is foundational. Cultural and governance misalignment is a leading cause of failed integrations.
	\end{block}
\end{frame}

\begin{frame}{D1-Q329: Regulatory Expectation for Governance}
	\begin{block}{Question}
		What do regulators MOST commonly expect regarding enterprise risk governance?
	\end{block}
	\begin{enumerate}[(A)]
		\item That all risk is eliminated
		\item That a documented, enforced governance framework exists
		\item That risk is managed only by external consultants
		\item That risk reporting is optional
	\end{enumerate}
\end{frame}
\begin{frame}{D1-Q329: Answer}
	\begin{block}{Correct Answer: B}
		Regulators expect documented governance frameworks with clear roles, accountability, and evidence of enforcement—not risk elimination.
	\end{block}
\end{frame}

\begin{frame}{D1-Q330: Risk Culture and Whistleblowing}
	\begin{block}{Question}
		Which mechanism BEST supports a healthy risk culture?
	\end{block}
	\begin{enumerate}[(A)]
		\item Anonymous whistleblowing channels with non-retaliation policies
		\item Strict punishment for all reported issues
		\item Discouraging escalation to avoid panic
		\item Restricting risk discussions to executives only
	\end{enumerate}
\end{frame}
\begin{frame}{D1-Q330: Answer}
	\begin{block}{Correct Answer: A}
		Anonymous reporting channels with strong non-retaliation protections encourage early disclosure and reduce the risk of hidden issues escalating.
	\end{block}
\end{frame}

\begin{frame}{D1-Q331: Governance of Third Parties}
	\begin{block}{Question}
		Which is MOST important for governing third-party risk at the enterprise level?
	\end{block}
	\begin{enumerate}[(A)]
		\item Delegating all oversight to the vendor
		\item A centralized third-party risk management program with clear accountability
		\item Reviewing vendor contracts annually only
		\item Relying on vendor self-assessments alone
	\end{enumerate}
\end{frame}
\begin{frame}{D1-Q331: Answer}
	\begin{block}{Correct Answer: B}
		A centralized program with clear accountability ensures consistent due diligence, contracting, monitoring, and termination practices across all third parties.
	\end{block}
\end{frame}

\begin{frame}{D1-Q332: Risk Governance and Business Strategy}
	\begin{block}{Question}
		How should risk governance relate to business strategy?
	\end{block}
	\begin{enumerate}[(A)]
		\item Risk governance should be independent and unrelated to strategy
		\item Risk governance should enable strategy by identifying and managing risk to objectives
		\item Risk governance should prevent any strategic change
		\item Risk governance should replace strategic planning
	\end{enumerate}
\end{frame}
\begin{frame}{D1-Q332: Answer}
	\begin{block}{Correct Answer: B}
		Effective risk governance enables strategy by ensuring risks to objectives are identified, assessed, and managed within appetite—not by blocking change.
	\end{block}
\end{frame}

\begin{frame}{D1-Q333: Risk Governance and Performance}
	\begin{block}{Question}
		Which is the BEST measure of risk governance effectiveness?
	\end{block}
	\begin{enumerate}[(A)]
		\item The number of risk policies published
		\item The degree to which risk-informed decisions are made and evidenced
		\item The size of the risk management department
		\item The frequency of risk committee meetings
	\end{enumerate}
\end{frame}
\begin{frame}{D1-Q333: Answer}
	\begin{block}{Correct Answer: B}
		Effectiveness is measured by whether risk information actually informs decisions—not by documentation volume, headcount, or meeting frequency.
	\end{block}
\end{frame}

\begin{frame}{D1-Q334: Governance and Risk Aggregation}
	\begin{block}{Question}
		Why is risk aggregation important for enterprise governance?
	\end{block}
	\begin{enumerate}[(A)]
		\item It eliminates the need for business-unit risk registers
		\item It provides a holistic view of exposures that may exceed appetite
		\item It ensures every risk has a unique owner
		\item It replaces the need for risk assessment
	\end{enumerate}
\end{frame}
\begin{frame}{D1-Q334: Answer}
	\begin{block}{Correct Answer: B}
		Aggregation reveals concentrations and cumulative exposures that may individually appear acceptable but collectively exceed appetite—critical for board oversight.
	\end{block}
\end{frame}

\begin{frame}{D1-Q335: Risk Governance in Acquisitions}
	\begin{block}{Question}
		During due diligence, which governance artifact is MOST valuable to review?
	\end{block}
	\begin{enumerate}[(A)]
		\item The target's marketing materials
		\item The target's risk register, appetite statement, and governance charters
		\item The target's employee newsletter
		\item The target's social media presence
	\end{enumerate}
\end{frame}
\begin{frame}{D1-Q335: Answer}
	\begin{block}{Correct Answer: B}
		Risk registers, appetite statements, and governance charters reveal the target's risk maturity, exposures, and accountability structures—essential for integration planning.
	\end{block}
\end{frame}

\begin{frame}{D1-Q336: Governance and Risk Appetite Revision}
	\begin{block}{Question}
		Who should approve changes to the enterprise risk appetite?
	\end{block}
	\begin{enumerate}[(A)]
		\item The risk management department
		\item The board of directors with senior management input
		\item Individual business unit managers
		\item The internal audit function
	\end{enumerate}
\end{frame}
\begin{frame}{D1-Q336: Answer}
	\begin{block}{Correct Answer: B}
		Risk appetite is a board-level decision informed by senior management. It reflects the enterprise's strategic direction and tolerance for risk.
	\end{block}
\end{frame}

\begin{frame}{D1-Q337: Risk Governance and Regulation}
	\begin{block}{Question}
		Which is the BEST way to integrate regulatory requirements into risk governance?
	\end{block}
	\begin{enumerate}[(A)]
		\item Treat compliance as a separate function from risk management
		\item Map regulatory obligations to risk categories and controls within the ERM framework
		\item Wait for regulators to identify gaps
		\item Delegate compliance entirely to external counsel
	\end{enumerate}
\end{frame}
\begin{frame}{D1-Q337: Answer}
	\begin{block}{Correct Answer: B}
		Mapping obligations to risk categories and controls ensures compliance is managed proactively within the ERM framework rather than as an isolated activity.
	\end{block}
\end{frame}

\begin{frame}{D1-Q338: Governance and Risk Tolerance}
	\begin{block}{Question}
		Which statement about risk tolerance is TRUE?
	\end{block}
	\begin{enumerate}[(A)]
		\item Tolerance is the same for every risk category
		\item Tolerance translates appetite into measurable limits for specific risks
		\item Tolerance is set by the internal audit function
		\item Tolerance is optional if appetite is documented
	\end{enumerate}
\end{frame}
\begin{frame}{D1-Q338: Answer}
	\begin{block}{Correct Answer: B}
		Tolerance operationalizes appetite by setting specific, measurable limits for individual risk categories, enabling monitoring and escalation.
	\end{block}
\end{frame}

\begin{frame}{D1-Q339: Governance and Culture Metrics}
	\begin{block}{Question}
		Which metric BEST reflects the health of an enterprise's risk culture?
	\end{block}
	\begin{enumerate}[(A)]
		\item Number of risk policies published
		\item Percentage of employees who report they feel safe raising risk concerns
		\item Number of risk committee meetings held
		\item Total spend on risk management tools
	\end{enumerate}
\end{frame}
\begin{frame}{D1-Q339: Answer}
	\begin{block}{Correct Answer: B}
		Psychological safety—measured through employee surveys—is a leading indicator of a healthy risk culture where issues are surfaced early.
	\end{block}
\end{frame}

\begin{frame}{D1-Q340: Governance and Emerging Risk}
	\begin{block}{Question}
		How should emerging risk be governed at the enterprise level?
	\end{block}
	\begin{enumerate}[(A)]
		\item Ignore emerging risk until it becomes a realized loss
		\item Establish a process for identifying, assessing, and escalating emerging risk
		\item Delegate emerging risk entirely to IT
		\item Report emerging risk only to external auditors
	\end{enumerate}
\end{frame}
\begin{frame}{D1-Q340: Answer}
	\begin{block}{Correct Answer: B}
		A structured process for identifying, assessing, and escalating emerging risk enables proactive management before issues crystallize into material losses.
	\end{block}
\end{frame}

\section{Domain 2: Additional Risk Assessment}

\begin{frame}{D2-Q341: Risk Assessment Scope}
	\begin{block}{Question}
		What is the FIRST step in defining the scope of a risk assessment?
	\end{block}
	\begin{enumerate}[(A)]
		\item Identify all possible threats
		\item Define the business objectives and processes in scope
		\item Select the risk assessment tool
		\item Assign risk owners
	\end{enumerate}
\end{frame}
\begin{frame}{D2-Q341: Answer}
	\begin{block}{Correct Answer: B}
		Scope is defined by the business objectives and processes under consideration. Threats, tools, and owners follow from a clear scope.
	\end{block}
\end{frame}

\begin{frame}{D2-Q342: Asset Valuation Methods}
	\begin{block}{Question}
		Which method is MOST appropriate for valuing an asset used in risk analysis?
	\end{block}
	\begin{enumerate}[(A)]
		\item Original purchase price only
		\item Replacement cost combined with business criticality
		\item Depreciation book value
		\item Insurance deductible amount
	\end{enumerate}
\end{frame}
\begin{frame}{D2-Q342: Answer}
	\begin{block}{Correct Answer: B}
		Asset value in risk analysis must reflect both replacement cost and the business impact of loss, including criticality to operations.
	\end{block}
\end{frame}

\begin{frame}{D2-Q343: Threat Modeling Purpose}
	\begin{block}{Question}
		What is the PRIMARY purpose of threat modeling in risk assessment?
	\end{block}
	\begin{enumerate}[(A)]
		\item To document all possible attack vectors
		\item To understand how threats could exploit vulnerabilities to achieve objectives
		\item To replace the need for vulnerability scanning
		\item To assign blame for past incidents
	\end{enumerate}
\end{frame}
\begin{frame}{D2-Q343: Answer}
	\begin{block}{Correct Answer: B}
		Threat modeling systematically examines how potential threats could exploit weaknesses to achieve adversarial or accidental objectives, informing risk prioritization.
	\end{block}
\end{frame}

\begin{frame}{D2-Q344: Qualitative Risk Scoring}
	\begin{block}{Question}
		Which is a limitation of qualitative risk scoring?
	\end{block}
	\begin{enumerate}[(A)]
		\item It is always more expensive than quantitative analysis
		\item Scores can be subjective and inconsistent across assessors
		\item It requires specialized statistical expertise
		\item It cannot be used for prioritization
	\end{enumerate}
\end{frame}
\begin{frame}{D2-Q344: Answer}
	\begin{block}{Correct Answer: B}
		Qualitative scoring depends on expert judgment and can vary between assessors. Calibration and clear criteria help mitigate subjectivity.
	\end{block}
\end{frame}

\begin{frame}{D2-Q345: Risk Scenario Development}
	\begin{block}{Question}
		Which is the BEST source of input for developing risk scenarios?
	\end{block}
	\begin{enumerate}[(A)]
		\item Vendor marketing materials
		\item Historical incidents, threat intelligence, and business objectives
		\item Social media trends
		\item Only internal audit findings
	\end{enumerate}
\end{frame}
\begin{frame}{D2-Q345: Answer}
	\begin{block}{Correct Answer: B}
		Effective risk scenarios draw on historical incidents, current threat intelligence, and business objectives to ensure relevance and realism.
	\end{block}
\end{frame}

\begin{frame}{D2-Q346: Risk Heat Map}
	\begin{block}{Question}
		What is the PRIMARY purpose of a risk heat map?
	\end{block}
	\begin{enumerate}[(A)]
		\item To calculate exact loss amounts
		\item To visualize and communicate risk based on likelihood and impact
		\item To replace the risk register
		\item To document control procedures
	\end{enumerate}
\end{frame}
\begin{frame}{D2-Q346: Answer}
	\begin{block}{Correct Answer: B}
		A heat map displays risks by likelihood and impact, enabling quick visual comparison and prioritization for stakeholders.
	\end{block}
\end{frame}

\begin{frame}{D2-Q347: Risk Register Updates}
	\begin{block}{Question}
		When should the risk register be updated?
	\end{block}
	\begin{enumerate}[(A)]
		\item Only during the annual risk assessment
		\item Continuously as new information becomes available
		\item Only when incidents occur
		\item Only when regulators request it
	\end{enumerate}
\end{frame}
\begin{frame}{D2-Q347: Answer}
	\begin{block}{Correct Answer: B}
		The risk register should be a living document updated as new risks emerge, controls change, or risk levels shift—not only on an annual cycle.
	\end{block}
\end{frame}

\begin{frame}{D2-Q348: Inherent Risk Assessment}
	\begin{block}{Question}
		What does inherent risk assessment measure?
	\end{block}
	\begin{enumerate}[(A)]
		\item Risk after all controls are implemented
		\item Risk before considering any controls
		\item Risk related only to financial loss
		\item Risk accepted by management
	\end{enumerate}
\end{frame}
\begin{frame}{D2-Q348: Answer}
	\begin{block}{Correct Answer: B}
		Inherent risk is the level of risk present before any controls are considered, providing a baseline for evaluating control effectiveness.
	\end{block}
\end{frame}

\begin{frame}{D2-Q349: Residual Risk Monitoring}
	\begin{block}{Question}
		Why must residual risk be monitored continuously?
	\end{block}
	\begin{enumerate}[(A)]
		\item Because controls degrade or become outdated over time
		\item Because inherent risk always decreases
		\item Because threats never change
		\item Because regulatory requirements are static
	\end{enumerate}
\end{frame}
\begin{frame}{D2-Q349: Answer}
	\begin{block}{Correct Answer: A}
		Controls can degrade, become misconfigured, or be bypassed, and threats and business conditions evolve—so residual risk must be continually reassessed.
	\end{block}
\end{frame}

\begin{frame}{D2-Q350: Risk Tolerance and Assessment}
	\begin{block}{Question}
		How does risk tolerance affect the risk assessment process?
	\end{block}
	\begin{enumerate}[(A)]
		\item It determines which risks require treatment and escalation
		\item It eliminates the need for risk assessment
		\item It makes all risks acceptable
		\item It replaces the risk register
	\end{enumerate}
\end{frame}
\begin{frame}{D2-Q350: Answer}
	\begin{block}{Correct Answer: A}
		Risk tolerance defines thresholds above which risks must be treated or escalated, guiding prioritization and response within the assessment process.
	\end{block}
\end{frame}

\begin{frame}{D2-Q351: Risk Assessment Frequency}
	\begin{block}{Question}
		What is the BEST approach for determining how often risk assessments should be performed?
	\end{block}
	\begin{enumerate}[(A)]
		\item Always annually regardless of changes
		\item Based on the rate of change in the risk environment and criticality of assets
		\item Only when mandated by regulation
		\item Never, once the initial assessment is complete
	\end{enumerate}
\end{frame}
\begin{frame}{D2-Q351: Answer}
	\begin{block}{Correct Answer: B}
		Assessment frequency should reflect the pace of change in threats, vulnerabilities, and business conditions, and the criticality of assets assessed.
	\end{block}
\end{frame}

\begin{frame}{D2-Q352: Risk Appetite and Risk Assessment}
	\begin{block}{Question}
		How does risk appetite inform risk assessment results?
	\end{block}
	\begin{enumerate}[(A)]
		\item It determines which risks are within acceptable bounds and which require action
		\item It eliminates the need for risk scoring
		\item It applies only to financial risks
		\item It prevents any risk from being identified
	\end{enumerate}
\end{frame}
\begin{frame}{D2-Q352: Answer}
	\begin{block}{Correct Answer: A}
		Risk appetite provides the reference point against which assessed risk levels are compared to determine acceptability and the need for treatment.
	\end{block}
\end{frame}

\begin{frame}{D2-Q353: Monte Carlo Simulation}
	\begin{block}{Question}
		When is Monte Carlo simulation MOST useful in risk assessment?
	\end{block}
	\begin{enumerate}[(A)]
		\item When a single deterministic outcome is needed
		\item When modeling the range of possible outcomes from multiple uncertain inputs
		\item When qualitative judgment alone is sufficient
		\item When only one variable affects the outcome
	\end{enumerate}
\end{frame}
\begin{frame}{D2-Q353: Answer}
	\begin{block}{Correct Answer: B}
		Monte Carlo simulation models the distribution of possible outcomes given multiple uncertain inputs, providing a probabilistic view of risk exposure.
	\end{block}
\end{frame}

\begin{frame}{D2-Q354: Risk Register Completeness}
	\begin{block}{Question}
		Which is the MOST important criterion for a complete risk register?
	\end{block}
	\begin{enumerate}[(A)]
		\item Every conceivable risk, no matter how trivial, is listed
		\item All material risks to business objectives are identified and assessed
		\item Only technical risks are documented
		\item Only risks with known solutions are included
	\end{enumerate}
\end{frame}
\begin{frame}{D2-Q354: Answer}
	\begin{block}{Correct Answer: B}
		A complete register captures material risks to business objectives. Including every trivial risk dilutes focus; excluding material risks leaves exposure unmanaged.
	\end{block}
\end{frame}

\begin{frame}{D2-Q355: Risk Assessment Stakeholders}
	\begin{block}{Question}
		Who should be involved in a risk assessment to ensure it is comprehensive?
	\end{block}
	\begin{enumerate}[(A)]
		\item Only the IT security team
		\item Business process owners, IT, risk management, and compliance
		\item Only external auditors
		\item Only senior executives
	\end{enumerate}
\end{frame}
\begin{frame}{D2-Q355: Answer}
	\begin{block}{Correct Answer: B}
		Comprehensive risk assessment requires cross-functional input to capture business, technical, regulatory, and operational perspectives.
	\end{block}
\end{frame}

\begin{frame}{D2-Q356: Risk Assessment Reporting}
	\begin{block}{Question}
		Which is the MOST important element of a risk assessment report to management?
	\end{block}
	\begin{enumerate}[(A)]
		\item Detailed technical vulnerability listings
		\item Clear linkage between risks, business impact, and recommended actions
		\item Exhaustive threat catalogs
		\item Methodological appendices
	\end{enumerate}
\end{frame}
\begin{frame}{D2-Q356: Answer}
	\begin{block}{Correct Answer: B}
		Management needs to understand how risks affect business objectives and what actions are recommended—technical detail supports but does not replace this.
	\end{block}
\end{frame}

\begin{frame}{D2-Q357: Risk Assessment and Control Testing}
	\begin{block}{Question}
		How does control testing relate to risk assessment?
	\end{block}
	\begin{enumerate}[(A)]
		\item Control testing replaces risk assessment
		\item Control testing provides evidence that controls are reducing risk as expected
		\item Control testing is unrelated to risk
		\item Control testing only applies to financial controls
	\end{enumerate}
\end{frame}
\begin{frame}{D2-Q357: Answer}
	\begin{block}{Correct Answer: B}
		Control testing verifies that controls operate effectively, validating the residual risk estimates produced by the risk assessment.
	\end{block}
\end{frame}

\begin{frame}{D2-Q358: Risk Assessment for Cloud}
	\begin{block}{Question}
		Which additional consideration is MOST important in a cloud risk assessment?
	\end{block}
	\begin{enumerate}[(A)]
		\item Physical security of the provider's data center
		\item Shared responsibility model and provider assurance
		\item The provider's marketing claims
		\item The provider's stock price
	\end{enumerate}
\end{frame}
\begin{frame}{D2-Q358: Answer}
	\begin{block}{Correct Answer: B}
		Cloud risk assessment must clarify which controls the provider owns versus the customer under the shared responsibility model, supported by independent assurance.
	\end{block}
\end{frame}

\begin{frame}{D2-Q359: Risk Assessment and Business Impact}
	\begin{block}{Question}
		How should business impact analysis (BIA) and risk assessment be integrated?
	\end{block}
	\begin{enumerate}[(A)]
		\item BIA replaces risk assessment
		\item BIA provides business impact data that informs risk prioritization and treatment
		\item They should never be linked
		\item Risk assessment provides recovery time objectives
	\end{enumerate}
\end{frame}
\begin{frame}{D2-Q359: Answer}
	\begin{block}{Correct Answer: B}
		BIA quantifies business impact of disruptions, which risk assessment uses to prioritize risks and justify treatment decisions.
	\end{block}
\end{frame}

\begin{frame}{D2-Q360: Risk Assessment and Regulatory Requirements}
	\begin{block}{Question}
		How should regulatory requirements be treated in a risk assessment?
	\end{block}
	\begin{enumerate}[(A)]
		\item Ignored because they are not business risks
		\item Assessed as sources of risk and compliance obligations
		\item Assumed to be satisfied by existing controls
		\item Delegated entirely to legal counsel
	\end{enumerate}
\end{frame}
\begin{frame}{D2-Q360: Answer}
	\begin{block}{Correct Answer: B}
		Regulatory obligations create risk of noncompliance and should be assessed alongside other risks, with controls evaluated for adequacy.
	\end{block}
\end{frame}

\begin{frame}{D2-Q361: Risk Assessment Scenario}
	\begin{block}{Question}
		Which is the BEST example of a well-formed risk scenario?
	\end{block}
	\begin{enumerate}[(A)]
		\item "The system might fail"
		\item "A phishing attack could compromise credentials, leading to unauthorized access to customer data and regulatory penalties"
		\item "Hackers are bad"
		\item "We need better security"
	\end{enumerate}
\end{frame}
\begin{frame}{D2-Q361: Answer}
	\begin{block}{Correct Answer: B}
		A well-formed scenario describes cause, event, and business consequence, enabling meaningful likelihood and impact assessment.
	\end{block}
\end{frame}

\begin{frame}{D2-Q362: Risk Appetite and Risk Assessment Results}
	\begin{block}{Question}
		What should happen when risk assessment results exceed risk appetite?
	\end{block}
	\begin{enumerate}[(A)]
		\item Ignore the results
		\item Escalate to risk owners and determine treatment options
		\item Revise the assessment methodology
		\item Accept the risk without documentation
	\end{enumerate}
\end{frame}
\begin{frame}{D2-Q362: Answer}
	\begin{block}{Correct Answer: B}
		When risk exceeds appetite, it must be escalated to risk owners for treatment decisions—acceptance requires documented approval and compensating controls.
	\end{block}
\end{frame}

\begin{frame}{D2-Q363: Risk Assessment and Key Risk Indicators}
	\begin{block}{Question}
		How do KRIs relate to risk assessment?
	\end{block}
	\begin{enumerate}[(A)]
		\item KRIs replace the need for periodic risk assessment
		\item KRIs provide early warning that a risk assessment update may be needed
		\item KRIs are only used after incidents
		\item KRIs are unrelated to risk assessment
	\end{enumerate}
\end{frame}
\begin{frame}{D2-Q363: Answer}
	\begin{block}{Correct Answer: B}
		KRIs monitor changes in risk exposure over time. Breaching a KRI threshold signals that the risk assessment should be revisited.
	\end{block}
\end{frame}

\begin{frame}{D2-Q364: Risk Assessment and Emerging Risk}
	\begin{block}{Question}
		How should emerging risk be handled in the risk assessment process?
	\end{block}
	\begin{enumerate}[(A)]
		\item Wait until it becomes a realized risk
		\item Include a horizon-scanning process to identify and assess emerging risk
		\item Exclude emerging risk from the register
		\item Delegate entirely to IT
	\end{enumerate}
\end{frame}
\begin{frame}{D2-Q364: Answer}
	\begin{block}{Correct Answer: B}
		Horizon scanning identifies emerging risks early, allowing proactive assessment and treatment before they materialize.
	\end{block}
\end{frame}

\begin{frame}{D2-Q365: Risk Assessment and Business Objectives}
	\begin{block}{Question}
		Why must risk assessment be tied to business objectives?
	\end{block}
	\begin{enumerate}[(A)]
		\item To make assessments longer
		\item To ensure risk evaluation reflects what the business is trying to achieve
		\item To satisfy auditors only
		\item To avoid identifying technical risks
	\end{enumerate}
\end{frame}
\begin{frame}{D2-Q365: Answer}
	\begin{block}{Correct Answer: B}
		Risk has meaning only in relation to objectives. Aligning assessment with business objectives ensures relevance and enables risk-informed decision-making.
	\end{block}
\end{frame}

% ============================================================
% DOMAIN 3 --- ADDITIONAL RISK RESPONSE AND REPORTING (Q366–Q395)
% ============================================================
\section{Domain 3: Additional Risk Response and Reporting}

\begin{frame}{D3-Q366: Risk Treatment Plan Components}
	\begin{block}{Question}
		Which is the MOST important component of a risk treatment plan?
	\end{block}
	\begin{enumerate}[(A)]
		\item The name of the auditor who approved it
		\item Specific actions, owners, and target dates for reducing risk
		\item The total budget allocated
		\item The tool used to document it
	\end{enumerate}
\end{frame}
\begin{frame}{D3-Q366: Answer}
	\begin{block}{Correct Answer: B}
		A treatment plan must specify actions, accountable owners, and target dates. Without these elements, the plan cannot be executed or monitored.
	\end{block}
\end{frame}

\begin{frame}{D3-Q367: Risk Response and Residual Risk}
	\begin{block}{Question}
		What is the PRIMARY goal of selecting a risk response?
	\end{block}
	\begin{enumerate}[(A)]
		\item Eliminate all residual risk
		\item Reduce residual risk to within risk tolerance
		\item Increase inherent risk
		\item Avoid documenting the response
	\end{enumerate}
\end{frame}
\begin{frame}{D3-Q367: Answer}
	\begin{block}{Correct Answer: B}
		The goal is to reduce residual risk to acceptable levels within appetite and tolerance—not to eliminate all risk, which is usually impossible.
	\end{block}
\end{frame}

\begin{frame}{D3-Q368: Risk Response and Cost-Benefit}
	\begin{block}{Question}
		When should a risk response NOT be implemented?
	\end{block}
	\begin{enumerate}[(A)]
		\item When it reduces risk below tolerance
		\item When the cost exceeds the benefit of risk reduction
		\item When it is required by regulation
		\item When management approves it
	\end{enumerate}
\end{frame}
\begin{frame}{D3-Q368: Answer}
	\begin{block}{Correct Answer: B}
		If the cost of a control exceeds its benefit, the response is not economically justified. Regulatory requirements may override cost-benefit considerations.
	\end{block}
\end{frame}

\begin{frame}{D3-Q369: Risk Transfer Limitations}
	\begin{block}{Question}
		Which is a limitation of risk transfer?
	\end{block}
	\begin{enumerate}[(A)]
		\item It eliminates all risk
		\item It does not eliminate reputational or legal risk
		\item It is always more expensive than mitigation
		\item It is prohibited by regulators
	\end{enumerate}
\end{frame}
\begin{frame}{D3-Q369: Answer}
	\begin{block}{Correct Answer: B}
		Risk transfer (e.g., insurance) shifts financial impact but not reputational, legal, or operational consequences. The enterprise remains exposed to those.
	\end{block}
\end{frame}

\begin{frame}{D3-Q370: Risk Acceptance Criteria}
	\begin{block}{Question}
		Which is the BEST criterion for accepting a risk?
	\end{block}
	\begin{enumerate}[(A)]
		\item The risk is within the enterprise's risk appetite and tolerance
		\item The risk has never occurred before
		\item No one has complained about the risk
		\item Insurance is too expensive
	\end{enumerate}
\end{frame}
\begin{frame}{D3-Q370: Answer}
	\begin{block}{Correct Answer: A}
		Risk acceptance is appropriate when the risk falls within the enterprise's appetite and tolerance. Acceptance must be documented and approved by the risk owner.
	\end{block}
\end{frame}

\begin{frame}{D3-Q371: Risk Avoidance Consequences}
	\begin{block}{Question}
		What is a common consequence of risk avoidance?
	\end{block}
	\begin{enumerate}[(A)]
		\item Loss of business opportunity
		\item Increased residual risk
		\item Decreased compliance
		\item Elimination of inherent risk only
	\end{enumerate}
\end{frame}
\begin{frame}{D3-Q371: Answer}
	\begin{block}{Correct Answer: A}
		Avoidance eliminates risk by exiting the activity but often forfeits associated opportunities, making it unsuitable when the business value is high.
	\end{block}
\end{frame}

\begin{frame}{D3-Q372: Risk Mitigation Cost}
	\begin{block}{Question}
		What is the BEST way to justify the cost of a risk mitigation control?
	\end{block}
	\begin{enumerate}[(A)]
		\item Compare the cost with the expected reduction in risk exposure
		\item Present the cost alone
		\item Compare the cost with the enterprise's total IT budget
		\item Avoid discussing cost
	\end{enumerate}
\end{frame}
\begin{frame}{D3-Q372: Answer}
	\begin{block}{Correct Answer: A}
		Justification requires demonstrating that the control reduces risk by an amount whose value exceeds the control's total cost of ownership.
	\end{block}
\end{frame}

\begin{frame}{D3-Q373: Control Ownership}
	\begin{block}{Question}
		What is the PRIMARY responsibility of a control owner?
	\end{block}
	\begin{enumerate}[(A)]
		\item Approve the enterprise risk appetite
		\item Ensure the control operates effectively and is monitored
		\item Conduct independent audits
		\item Set the risk tolerance
	\end{enumerate}
\end{frame}
\begin{frame}{D3-Q373: Answer}
	\begin{block}{Correct Answer: B}
		Control owners are accountable for ensuring controls operate as designed, are monitored, and are updated as conditions change.
	\end{block}
\end{frame}

\begin{frame}{D3-Q374 Risk Response for Regulatory Risk}
	\begin{block}{Question}
		Which response is MOST appropriate for a regulatory compliance risk?
	\end{block}
	\begin{enumerate}[(A)]
		\item Accept the risk and hope for no enforcement
		\item Mitigate by implementing required controls and monitoring compliance
		\item Transfer the risk through insurance only
		\item Ignore the risk until inspected
	\end{enumerate}
\end{frame}
\begin{frame}{D3-Q374 Answer}
	\begin{block}{Correct Answer: B}
		Regulatory risk is typically mitigated through required controls and compliance monitoring. Acceptance is rarely justified given enforcement and reputational consequences.
	\end{block}
\end{frame}

\begin{frame}{D3-Q375: Risk Response and Stakeholder Engagement}
	\begin{block}{Question}
		Which stakeholders are MOST important to engage when selecting a risk response?
	\end{block}
	\begin{enumerate}[(A)]
		\item Only the IT department
		\item Business owners, risk owners, and affected operational teams
		\item Only external auditors
		\item Only the legal department
	\end{enumerate}
\end{frame}
\begin{frame}{D3-Q375: Answer}
	\begin{block}{Correct Answer: B}
		Those who own the risk and are affected by the response must be engaged to ensure the response is feasible, acceptable, and sustainable.
	\end{block}
\end{frame}

\begin{frame}{D3-Q376: Risk Response Monitoring}
	\begin{block}{Question}
		How should the effectiveness of a risk response be monitored?
	\end{block}
	\begin{enumerate}[(A)]
		\item Only through annual audits
		\item Through KPIs, KCIs, and periodic reviews of residual risk
		\item By asking the vendor
		\item By reviewing social media
	\end{enumerate}
\end{frame}
\begin{frame}{D3-Q376: Answer}
	\begin{block}{Correct Answer: B}
		Monitoring should combine performance indicators, control indicators, and periodic reassessment of residual risk against appetite.
	\end{block}
\end{frame}

\begin{frame}{D3-Q377: Risk Response Documentation}
	\begin{block}{Question}
		Why must risk responses be documented?
	\end{block}
	\begin{enumerate}[(A)]
		\item To increase paperwork
		\item To provide evidence of decisions, accountability, and monitoring
		\item To satisfy vendor requirements
		\item To avoid escalation
	\end{enumerate}
\end{frame}
\begin{frame}{D3-Q377: Answer}
	\begin{block}{Correct Answer: B}
		Documentation provides evidence that responses were chosen deliberately, owners assigned, and results monitored—essential for accountability and audit.
	\end{block}
\end{frame}

\begin{frame}{D3-Q378: Risk Response and Residual Risk Acceptance}
	\begin{block}{Question}
		Who should formally accept residual risk after treatment?
	\end{block}
	\begin{enumerate}[(A)]
		\item The control owner
		\item The risk owner with appropriate management approval
		\item The internal auditor
		\item The IT administrator
	\end{enumerate}
\end{frame}
\begin{frame}{D3-Q378: Answer}
	\begin{block}{Correct Answer: B}
		Residual risk acceptance rests with the risk owner, with approval from management at the appropriate level based on the risk's significance.
	\end{block}
\end{frame}

\begin{frame}{D3-Q379: Risk Response and Insurance}
	\begin{block}{Question}
		Which type of risk is MOST suitable for transfer through insurance?
	\end{block}
	\begin{enumerate}[(A)]
		\item High-likelihood, low-impact operational errors
		\item Low-likelihood, high-impact catastrophic events
		\item Risks that are entirely within the organization's control
		\item Risks with no measurable financial consequence
	\end{enumerate}
\end{frame}
\begin{frame}{D3-Q379: Answer}
	\begin{block}{Correct Answer: B}
		Insurance is most economically efficient for low-frequency, high-severity events that would otherwise threaten the enterprise's financial stability.
	\end{block}
\end{frame}

\begin{frame}{D3-Q380: Risk Reporting Frequency}
	\begin{block}{Question}
		What is the BEST approach for determining risk reporting frequency to senior management?
	\end{block}
	\begin{enumerate}[(A)]
		\item Always monthly
		\item Aligned with the pace of change and criticality of the risks reported
		\item Only annually
		\item Only after incidents
	\end{enumerate}
\end{frame}
\begin{frame}{D3-Q380: Answer}
	\begin{block}{Correct Answer: B}
		Reporting frequency should reflect how rapidly the risk environment changes and how critical those risks are to business objectives.
	\end{block}
\end{frame}

\begin{frame}{D3-Q381: Risk Reporting Content}
	\begin{block}{Question}
		Which content is MOST valuable in a risk report to senior management?
	\end{block}
	\begin{enumerate}[(A)]
		\item Raw vulnerability scan output
		\item Aggregated risk against appetite, key changes, and recommended actions
		\item Exhaustive threat intelligence feeds
		\item Detailed configuration settings
	\end{enumerate}
\end{frame}
\begin{frame}{D3-Q381: Answer}
	\begin{block}{Correct Answer: B}
		Executive reporting should focus on aggregated risk, changes since the last report, and recommended actions—not technical detail.
	\end{block}
\end{frame}

\begin{frame}{D3-Q382: Risk Response and Regulations}
	\begin{block}{Question}
		How should regulatory requirements influence risk response selection?
	\end{block}
	\begin{enumerate}[(A)]
		\item They can be ignored if mitigation is expensive
		\item They may mandate specific responses regardless of cost-benefit
		\item They are only relevant to financial institutions
		\item They should be delegated to external counsel
	\end{enumerate}
\end{frame}
\begin{frame}{D3-Q382: Answer}
	\begin{block}{Correct Answer: B}
		Regulatory requirements may compel specific controls or responses irrespective of cost-benefit analysis. Noncompliance carries penalties and reputational harm.
	\end{block}
\end{frame}

\begin{frame}{D3-Q383: Risk Response and Third Parties}
	\begin{block}{Question}
		Which is MOST important when relying on a third party for risk mitigation?
	\end{block}
	\begin{enumerate}[(A)]
		\item The third party's marketing claims
		\item Contractual requirements, audit rights, and monitoring
		\item The third party's size
		\item The third party's location
	\end{enumerate}
\end{frame}
\begin{frame}{D3-Q383: Answer}
	\begin{block}{Correct Answer: B}
		Contractual requirements, audit rights, and ongoing monitoring ensure the third party actually delivers the risk reduction expected.
	\end{block}
\end{frame}

\begin{frame}{D3-Q384: Risk Response and Residual Risk Target}
	\begin{block}{Question}
		What should be the target level of residual risk after treatment?
	\end{block}
	\begin{enumerate}[(A)]
		\item Zero
		\item Within the enterprise's risk appetite and tolerance
		\item As high as possible
		\item Undefined
	\end{enumerate}
\end{frame}
\begin{frame}{D3-Q384: Answer}
	\begin{block}{Correct Answer: B}
		The target is to bring residual risk within appetite and tolerance. Zero risk is generally neither achievable nor economically justified.
	\end{block}
\end{frame}

\begin{frame}{D3-Q385: Risk Response and Change Management}
	\begin{block}{Question}
		How does change management relate to risk response?
	\end{block}
	\begin{enumerate}[(A)]
		\item Change management is unrelated to risk
		\item Changes can introduce new risks and must be assessed before implementation
		\item Change management replaces risk response
		\item All changes are risk-free if approved
	\end{enumerate}
\end{frame}
\begin{frame}{D3-Q385: Answer}
	\begin{block}{Correct Answer: B}
		Changes to systems, processes, or controls can introduce new risks or alter existing ones, so they must be assessed through change management.
	\end{block}
\end{frame}

\begin{frame}{D3-Q386: Risk Response and Emerging Risk}
	\begin{block}{Question}
		How should risk responses address emerging risk?
	\end{block}
	\begin{enumerate}[(A)]
		\item Ignore emerging risk until it materializes
		\item Monitor emerging risk and prepare contingency responses
		\item Treat emerging risk the same as known risk without assessment
		\item Delegate to the vendor
	\end{enumerate}
\end{frame}
\begin{frame}{D3-Q386: Answer}
	\begin{block}{Correct Answer: B}
		Emerging risk should be monitored through horizon scanning, with contingency responses prepared for plausible scenarios.
	\end{block}
\end{frame}

\begin{frame}{D3-Q387: Risk Response and Business Continuity}
	\begin{block}{Question}
		How does business continuity planning relate to risk response?
	\end{block}
	\begin{enumerate}[(A)]
		\item It replaces risk response
		\item It provides a response for risks that cannot be fully mitigated
		\item It is unrelated to risk response
		\item It is only for natural disasters
	\end{enumerate}
\end{frame}
\begin{frame}{D3-Q387: Answer}
	\begin{block}{Correct Answer: B}
		Business continuity provides response and recovery capabilities for risks that cannot be fully mitigated, ensuring resilience.
	\end{block}
\end{frame}

\begin{frame}{D3-Q388: Risk Response and Control Design}
	\begin{block}{Question}
		What is the FIRST consideration when designing a control as a risk response?
	\end{block}
	\begin{enumerate}[(A)]
		\item Cost of the control
		\item The specific risk and its root cause the control is meant to address
		\item The vendor's reputation
		\item The control's popularity in the industry
	\end{enumerate}
\end{frame}
\begin{frame}{D3-Q388: Answer}
	\begin{block}{Correct Answer: B}
		Controls must be designed to address the specific risk and its root cause. Otherwise, the control may be ineffective regardless of cost or popularity.
	\end{block}
\end{frame}

\begin{frame}{D3-Q389: Risk Response and Monitoring Metrics}
	\begin{block}{Question}
		Which metric BEST indicates a risk response is effective?
	\end{block}
	\begin{enumerate}[(A)]
		\item Number of controls implemented
		\item Reduction in residual risk relative to tolerance
		\item Amount spent on controls
		\item Frequency of control audits
	\end{enumerate}
\end{frame}
\begin{frame}{D3-Q389: Answer}
	\begin{block}{Correct Answer: B}
		Effectiveness is measured by the reduction in residual risk achieved relative to tolerance—not by activity metrics.
	\end{block}
\end{frame}

\begin{frame}{D3-Q390: Risk Response and Residual Risk Escalation}
	\begin{block}{Question}
		When should residual risk be escalated?
	\end{block}
	\begin{enumerate}[(A)]
		\item Never
		\item When it exceeds tolerance or appetite
		\item Only after an incident
		\item Only if the auditor requests it
	\end{enumerate}
\end{frame}
\begin{frame}{D3-Q390: Answer}
	\begin{block}{Correct Answer: B}
		Residual risk exceeding tolerance or appetite must be escalated to appropriate management for decision on further treatment or acceptance.
	\end{block}
\end{frame}

\begin{frame}{D3-Q391: Risk Response and Accountability}
	\begin{block}{Question}
		Who is accountable for the effectiveness of a risk response?
	\end{block}
	\begin{enumerate}[(A)]
		\item The control vendor
		\item The risk owner
		\item The internal auditor
		\item The IT administrator
	\end{enumerate}
\end{frame}
\begin{frame}{D3-Q391: Answer}
	\begin{block}{Correct Answer: B}
		The risk owner is accountable for ensuring the chosen response is implemented, monitored, and effective.
	\end{block}
\end{frame}

\begin{frame}{D3-Q392: Risk Response and Documentation}
	\begin{block}{Question}
		Which element is MOST important to document for a risk response?
	\end{block}
	\begin{enumerate}[(A)]
		\item The response rationale, owner, and target date
		\item The vendor's contact information
		\item The control's technical specification
		\item The marketing brochure
	\end{enumerate}
\end{frame}
\begin{frame}{D3-Q392: Answer}
	\begin{block}{Correct Answer: A}
		Documenting the rationale, owner, and target date provides accountability and enables monitoring. Technical details may be additional but are not primary.
	\end{block}
\end{frame}

\begin{frame}{D3-Q393: Risk Response and Regulatory Reporting}
	\begin{block}{Question}
		How should regulatory reporting requirements affect risk response documentation?
	\end{block}
	\begin{enumerate}[(A)]
		\item They have no impact
		\item They may dictate specific documentation and reporting obligations
		\item They only apply to financial firms
		\item They can be ignored
	\end{enumerate}
\end{frame}
\begin{frame}{D3-Q393: Answer}
	\begin{block}{Correct Answer: B}
		Regulatory requirements may impose specific documentation, retention, and reporting obligations that must be reflected in risk response processes.
	\end{block}
\end{frame}

\begin{frame}{D3-Q394: Risk Response and Residual Risk Reporting}
	\begin{block}{Question}
		Which is MOST valuable in reporting residual risk to management?
	\end{block}
	\begin{enumerate}[(A)]
		\item A comparison of residual risk against appetite and tolerance with trends
		\item Raw control inventory
		\item The number of policies published
		\item Vendor performance metrics
	\end{enumerate}
\end{frame}
\begin{frame}{D3-Q394: Answer}
	\begin{block}{Correct Answer: A}
		Management needs to see residual risk relative to appetite and tolerance, with trends over time, to make informed decisions.
	\end{block}
\end{frame}

\begin{frame}{D3-Q395: Risk Response and Continuous Improvement}
	\begin{block}{Question}
		How should lessons learned from risk responses be used?
	\end{block}
	\begin{enumerate}[(A)]
		\item Filed and forgotten
		\item To improve future risk assessments and response selection
		\item To assign blame
		\item Shared only with auditors
	\end{enumerate}
\end{frame}
\begin{frame}{D3-Q395: Answer}
	\begin{block}{Correct Answer: B}
		Lessons learned should feed back into the risk management process to improve future assessments, response selection, and control design.
	\end{block}
\end{frame}


\section{Domain 4: Additional IT and Security}

\begin{frame}{D4-Q396: Security Architecture Purpose}
	\begin{block}{Question}
		What is the PRIMARY purpose of an enterprise security architecture?
	\end{block}
	\begin{enumerate}[(A)]
		\item To replace the risk register
		\item To align security controls and strategies with business objectives
		\item To document every vulnerability
		\item To eliminate the need for security policies
	\end{enumerate}
\end{frame}
\begin{frame}{D4-Q396: Answer}
	\begin{block}{Correct Answer: B}
		Enterprise security architecture aligns security strategies and controls with business objectives, ensuring coherent and integrated protection.
	\end{block}
\end{frame}

\begin{frame}{D4-Q397: Access Control Models}
	\begin{block}{Question}
		Which access control model is MOST appropriate for an enterprise with highly dynamic roles?
	\end{block}
	\begin{enumerate}[(A)]
		\item Mandatory access control (MAC)
		\item Role-based access control (RBAC)
		\item Discretionary access control (DAC)
		\item Rule-based access control without roles
	\end{enumerate}
\end{frame}
\begin{frame}{D4-Q397: Answer}
	\begin{block}{Correct Answer: B}
		RBAC maps access to roles rather than individuals, simplifying administration and accommodating dynamic role assignments.
	\end{block}
\end{frame}

\begin{frame}{D4-Q398: Encryption Key Management}
	\begin{block}{Question}
		Which is the MOST critical aspect of encryption key management?
	\end{block}
	\begin{enumerate}[(A)]
		\item Publishing keys for transparency
		\item Secure generation, distribution, storage, and rotation of keys
		\item Using the same key indefinitely
		\item Storing keys with the encrypted data
	\end{enumerate}
\end{frame}
\begin{frame}{D4-Q398: Answer}
	\begin{block}{Correct Answer: B}
		Key management must ensure secure generation, distribution, storage, rotation, and destruction of keys—otherwise encryption provides no protection.
	\end{block}
\end{frame}

\begin{frame}{D4-Q399: Security Awareness Training}
	\begin{block}{Question}
		What is the MOST effective format for ongoing security awareness training?
	\end{block}
	\begin{enumerate}[(A)]
		\item Annual classroom training only
		\item Continuous, role-specific, and reinforced through simulations
		\item A single email at hire
		\item A poster in the break room
	\end{enumerate}
\end{frame}
\begin{frame}{D4-Q399: Answer}
	\begin{block}{Correct Answer: B}
		Continuous, role-specific training reinforced by phishing simulations and other exercises keeps security awareness current and effective.
	\end{block}
\end{frame}

\begin{frame}{D4-Q400: Security Incident Response Phases}
	\begin{block}{Question}
		Which phase of incident response focuses on limiting the impact of an incident?
	\end{block}
	\begin{enumerate}[(A)]
		\item Preparation
		\item Detection and analysis
		\item Containment, eradication, and recovery
		\item Post-incident activity
	\end{enumerate}
\end{frame}
\begin{frame}{D4-Q400: Answer}
	\begin{block}{Correct Answer: C}
		Containment limits the spread and impact of an incident; eradication removes the cause; recovery restores operations. Post-incident activity reviews lessons learned.
	\end{block}
\end{frame}

\begin{frame}{D4-Q401: Vulnerability Management}
	\begin{block}{Question}
		Which is the MOST important step after identifying a critical vulnerability?
	\end{block}
	\begin{enumerate}[(A)]
		\item Document it and wait
		\item Assess its risk and prioritize remediation based on business impact
		\item Immediately shut down all systems
		\item Blame the vendor
	\end{enumerate}
\end{frame}
\begin{frame}{D4-Q401: Answer}
	\begin{block}{Correct Answer: B}
		Vulnerabilities must be assessed in the context of business impact and exploitability, then prioritized for remediation based on risk.
	\end{block}
\end{frame}

\begin{frame}{D4-Q402: Data Loss Prevention}
	\begin{block}{Question}
		What is the PRIMARY function of a data loss prevention (DLP) system?
	\end{block}
	\begin{enumerate}[(A)]
		\item Encrypt all data at rest
		\item Detect and prevent unauthorized transfer of sensitive data
		\item Back up all files
		\item Manage user passwords
	\end{enumerate}
\end{frame}
\begin{frame}{D4-Q402: Answer}
	\begin{block}{Correct Answer: B}
		DLP monitors, detects, and blocks unauthorized transmission of sensitive data based on policy rules and data classification.
	\end{block}
\end{frame}

\begin{frame}{D4-Q403: Identity and Access Management}
	\begin{block}{Question}
		Which principle is MOST fundamental to identity and access management?
	\end{block}
	\begin{enumerate}[(A)]
		\item Shared accounts for convenience
		\item Least privilege and segregation of duties
		\item Permanent access for all employees
		\item Access without authentication
	\end{enumerate}
\end{frame}
\begin{frame}{D4-Q403: Answer}
	\begin{block}{Correct Answer: B}
		Least privilege and segregation of duties minimize access and prevent conflicts of interest, forming the foundation of IAM.
	\end{block}
\end{frame}

\begin{frame}{D4-Q404: Security Monitoring}
	\begin{block}{Question}
		Which is the MOST important factor for effective security monitoring?
	\end{block}
	\begin{enumerate}[(A)]
		\item Buying the most expensive SIEM tool
		\item Trained staff who can interpret and act on alerts
		\item Collecting all possible logs without filtering
		\item Monitoring only during business hours
	\end{enumerate}
\end{frame}
\begin{frame}{D4-Q404: Answer}
	\begin{block}{Correct Answer: B}
		Tools are only as effective as the people operating them. Trained analysts who can triage, investigate, and respond are essential.
	\end{block}
\end{frame}

\begin{frame}{D4-Q405: Patch Management}
	\begin{block}{Question}
		What is the BEST practice for patch management?
	\end{block}
	\begin{enumerate}[(A)]
		\item Apply all patches immediately without testing
		\item Test patches in a controlled environment before production deployment
		\item Never apply patches to avoid disruption
		\item Delegate patching entirely to vendors
	\end{enumerate}
\end{frame}
\begin{frame}{D4-Q405: Answer}
	\begin{block}{Correct Answer: B}
		Patches should be tested in a controlled environment before production rollout to avoid introducing new issues, while still applying them promptly.
	\end{block}
\end{frame}

\begin{frame}{D4-Q406: Cloud Shared Responsibility}
	\begin{block}{Question}
		Under the cloud shared responsibility model, what is the customer ALWAYS responsible for?
	\end{block}
	\begin{enumerate}[(A)]
		\item Physical security of the data center
		\item Data classification, access management, and configuration of customer-controlled components
		\item Hypervisor patching
		\item Network cabling
	\end{enumerate}
\end{frame}
\begin{frame}{D4-Q406: Answer}
	\begin{block}{Correct Answer: B}
		Customers retain responsibility for their data, access controls, and configuration of components within their control—regardless of cloud model.
	\end{block}
\end{frame}

\begin{frame}{D4-Q407: Third-Party Risk Assessment}
	\begin{block}{Question}
		Which is the MOST important element of a third-party risk assessment?
	\end{block}
	\begin{enumerate}[(A)]
		\item The vendor's marketing materials
		\item Evidence of security control effectiveness (e.g., SOC 2, ISO 27001)
		\item The vendor's number of employees
		\item The vendor's office location
	\end{enumerate}
\end{frame}
\begin{frame}{D4-Q407: Answer}
	\begin{block}{Correct Answer: B}
		Independent evidence of control effectiveness such as SOC 2 reports or ISO 27001 certification provides reliable assurance.
	\end{block}
\end{frame}

\begin{frame}{D4-Q408: Data Retention and Disposal}
	\begin{block}{Question}
		What is the MOST important consideration for data disposal?
	\end{block}
	\begin{enumerate}[(A)]
		\item Delete files quickly to save space
		\item Ensure data is unrecoverable in accordance with policy and regulation
		\item Print all data before disposal
		\item Archive everything forever
	\end{enumerate}
\end{frame}
\begin{frame}{D4-Q408: Answer}
	\begin{block}{Correct Answer: B}
		Secure disposal must render data unrecoverable, consistent with retention policies and regulatory requirements.
	\end{block}
\end{frame}

\begin{frame}{D4-Q409: Application Security Testing}
	\begin{block}{Question}
		Which type of testing is MOST effective for finding security flaws in source code?
	\end{block}
	\begin{enumerate}[(A)]
		\item Penetration testing only
		\item Static application security testing (SAST)
		\item Vulnerability scanning of the network
		\item User acceptance testing
	\end{enumerate}
\end{frame}
\begin{frame}{D4-Q409: Answer}
	\begin{block}{Correct Answer: B}
		SAST analyzes source code for security weaknesses without executing the application, making it effective for finding code-level flaws.
	\end{block}
\end{frame}

\begin{frame}{D4-Q410: Endpoint Protection}
	\begin{block}{Question}
		Which is the MOST important element of endpoint protection in a modern enterprise?
	\end{block}
	\begin{enumerate}[(A)]
		\item Antivirus only
		\item Layered defenses including EDR, patching, and configuration management
		\item Firewall only
		\item Password complexity only
	\end{enumerate}
\end{frame}
\begin{frame}{D4-Q410: Answer}
	\begin{block}{Correct Answer: B}
		Effective endpoint protection requires layered defenses—EDR, patching, configuration management, and user awareness—not any single tool.
	\end{block}
\end{frame}

\begin{frame}{D4-Q411: Network Segmentation}
	\begin{block}{Question}
		What is the PRIMARY purpose of network segmentation?
	\end{block}
	\begin{enumerate}[(A)]
		\item To increase bandwidth
		\item To limit the spread of attacks and protect sensitive segments
		\item To simplify routing only
		\item To reduce cable costs
	\end{enumerate}
\end{frame}
\begin{frame}{D4-Q411: Answer}
	\begin{block}{Correct Answer: B}
		Segmentation isolates sensitive assets and limits lateral movement, reducing the impact of a breach.
	\end{block}
\end{frame}

\begin{frame}{D4-Q412: Log Management}
	\begin{block}{Question}
		Which is the MOST important consideration for log management?
	\end{block}
	\begin{enumerate}[(A)]
		\item Collecting logs and never reviewing them
		\item Centralized, protected storage with timely review and retention per policy
		\item Deleting logs immediately after generation
		\item Storing logs only on the source system
	\end{enumerate}
\end{frame}
\begin{frame}{D4-Q412: Answer}
	\begin{block}{Correct Answer: B}
		Logs must be centrally collected, protected from tampering, reviewed in a timely manner, and retained per policy to support detection and forensics.
	\end{block}
\end{frame}

\begin{frame}{D4-Q413: Security Metrics}
	\begin{block}{Question}
		Which is the BEST security metric for executive reporting?
	\end{block}
	\begin{enumerate}[(A)]
		\item Number of firewall rules
		\item Trends in critical vulnerabilities, incidents, and control effectiveness relative to tolerance
		\item Count of security policies
		\item Number of antivirus definitions
	\end{enumerate}
\end{frame}
\begin{frame}{D4-Q413: Answer}
	\begin{block}{Correct Answer: B}
		Executives need trend-based metrics on critical vulnerabilities, incidents, and control effectiveness relative to tolerance—not operational counts.
	\end{block}
\end{frame}

\begin{frame}{D4-Q414: Business Continuity Testing}
	\begin{block}{Question}
		What is the MOST effective way to validate a business continuity plan?
	\end{block}
	\begin{enumerate}[(A)]
		\item Reviewing the document only
		\item Exercising the plan through tests and simulations and updating based on results
		\item Filing it with regulators
		\item Sharing it with the team
	\end{enumerate}
\end{frame}
\begin{frame}{D4-Q414: Answer}
	\begin{block}{Correct Answer: B}
		Plans must be exercised—through tabletop exercises, simulations, or full tests—and updated based on lessons learned to remain effective.
	\end{block}
\end{frame}

\begin{frame}{D4-Q415: Security Culture}
	\begin{block}{Question}
		Which BEST describes the relationship between security culture and technical controls?
	\end{block}
	\begin{enumerate}[(A)]
		\item Technical controls alone are sufficient
		\item Security culture and technical controls are complementary and both essential
		\item Culture is irrelevant if controls are strong
		\item Culture replaces technical controls
	\end{enumerate}
\end{frame}
\begin{frame}{D4-Q415: Answer}
	\begin{block}{Correct Answer: B}
		Strong technical controls can be undermined by poor culture (e.g., password sharing), and good culture cannot compensate for missing controls. Both are essential.
	\end{block}
\end{frame}


	% ============================================================
% DOMAIN 1 --- GOVERNANCE: PRACTICAL AND SCENARIO QUESTIONS (Q516–Q540)
% ============================================================
\section{Domain 1: Governance Scenarios}

\begin{frame}{D1-Q516: Enterprise Risk Committee Charter}
	\begin{block}{Question}
		An enterprise is forming a risk committee. Which element is MOST important to include in its charter?
	\end{block}
	\begin{enumerate}[(A)]
		\item The names of all risk owners
		\item Its authority, scope, and reporting lines to the board
		\item The tools it will use
		\item Its meeting schedule
	\end{enumerate}
\end{frame}
\begin{frame}{D1-Q516: Answer}
	\begin{block}{Correct Answer: B}
		A charter must define authority, scope, and reporting relationships. Without these, the committee cannot provide effective oversight or enforce decisions.
	\end{block}
\end{frame}

\begin{frame}{D1-Q517: Business Unit Risk Reporting}
	\begin{block}{Question}
		A business unit routinely submits risk reports late. What is the MOST likely governance root cause?
	\end{block}
	\begin{enumerate}[(A)]
		\item The reporting tool is inefficient
		\item Unclear accountability and lack of consequences for late reporting
		\item Too few risk reports required
		\item Excessive regulation
	\end{enumerate}
\end{frame}
\begin{frame}{D1-Q517: Answer}
	\begin{block}{Correct Answer: B}
		Late reporting usually signals unclear accountability or missing consequences. Governance must define deadlines, owners, and escalation for noncompliance.
	\end{block}
\end{frame}

\begin{frame}{D1-Q518: Risk Appetite and Budget}
	\begin{block}{Question}
		Management proposes a major investment outside the enterprise's risk appetite. What should the risk practitioner do FIRST?
	\end{block}
	\begin{enumerate}[(A)]
		\item Approve the investment
		\item Escalate the misalignment to the risk owner and senior management
		\item Ignore because investments are business decisions
		\item Revise the risk appetite
	\end{enumerate}
\end{frame}
\begin{frame}{D1-Q518: Answer}
	\begin{block}{Correct Answer: B}
		The risk practitioner must escalate the misalignment. Approving or silently ignoring a decision outside appetite undermines governance.
	\end{block}
\end{frame}

\begin{frame}{D1-Q519: Risk Culture and KPIs}
	\begin{block}{Question}
		Which KPI BEST reflects the health of an enterprise risk culture?
	\end{block}
	\begin{enumerate}[(A)]
		\item Number of risk assessments completed
		\item Percentage of risk owners who completed their attestations on time
		\item Number of risk policies published
		\item Total spend on risk tools
	\end{enumerate}
\end{frame}
\begin{frame}{D1-Q519: Answer}
	\begin{block}{Correct Answer: B}
		Timely attestations indicate engaged risk owners and a functioning culture—not just completed paperwork.
	\end{block}
\end{frame}

\begin{frame}{D1-Q520: Board Members and Risk Detail}
	\begin{block}{Question}
		A board member requests a list of every open vulnerability. What is the BEST response?
	\end{block}
	\begin{enumerate}[(A)]
		\item Provide the full list
		\item Provide aggregated risk information aligned to appetite, with vulnerability detail available on request
		\item Refuse the request
		\item Delay until the annual report
	\end{enumerate}
\end{frame}
\begin{frame}{D1-Q520: Answer}
	\begin{block}{Correct Answer: B}
		Boards need aggregated risk against appetite. Raw vulnerability lists obscure judgment; aggregated views plus on-request detail support effective oversight.
	\end{block}
\end{frame}

\begin{frame}{D1-Q521: Risk Acceptance Authority}
	\begin{block}{Question}
		Who should have authority to accept enterprise-level risk?
	\end{block}
	\begin{enumerate}[(A)]
		\item The risk practitioner
		\item Business unit managers
		\item Senior management or board, within delegated limits
		\item External auditors
	\end{enumerate}
\end{frame}
\begin{frame}{D1-Q521: Answer}
	\begin{block}{Correct Answer: C}
		Risk acceptance authority must sit with management at the appropriate level, defined by limits and delegated formally.
	\end{block}
\end{frame}

\begin{frame}{D1-Q522: Emerging Risk Governance}
	\begin{block}{Question}
		Which is the BEST governance mechanism for emerging risk?
	\end{block}
	\begin{enumerate}[(A)]
		\item Ignore until it crystallizes
		\item Horizon scanning with periodic reporting to the risk committee
		\item Delegate entirely to IT
		\item Report only after incidents
	\end{enumerate}
\end{frame}
\begin{frame}{D1-Q522: Answer}
	\begin{block}{Correct Answer: B}
		Horizon scanning with periodic reporting ensures emerging risk is identified and assessed before it becomes a material exposure.
	\end{block}
\end{frame}

\begin{frame}{D1-Q523: Risk Awareness Program Failure}
	\begin{block}{Question}
		A risk awareness program fails to change behavior. What is the MOST likely cause?
	\end{block}
	\begin{enumerate}[(A)]
		\item The intranet site is unattractive
		\item Training is generic and not tied to employees' roles
		\item Not enough posters
		\item The program runs too often
	\end{enumerate}
\end{frame}
\begin{frame}{D1-Q523: Answer}
	\begin{block}{Correct Answer: B}
		Role-relevant training is essential for behavior change. Generic messages are ignored; tailored scenarios drive action.
	\end{block}
\end{frame}

\begin{frame}{D1-Q524: Risk Appetite Consistency}
	\begin{block}{Question}
		An enterprise's risk appetite statement conflicts with several operational policies. What should be done FIRST?
	\end{block}
	\begin{enumerate}[(A)]
		\item Change the appetite statement
		\item Align policies to appetite through review and revision
		\item Ignore the inconsistency
		\item Suspend operations
	\end{enumerate}
\end{frame}
\begin{frame}{D1-Q524: Answer}
	\begin{block}{Correct Answer: B}
		Policies must be consistent with appetite. Misalignment indicates a governance gap and can cause risk-taking outside tolerance.
	\end{block}
\end{frame}

\begin{frame}{D1-Q525: Third-Line Independence}
	\begin{block}{Question}
		Internal audit reports to the CRO rather than the board. What is the MOST serious consequence?
	\end{block}
	\begin{enumerate}[(A)]
		\item Increased workload for audit
		\item Reduced independence and objectivity of assurance
		\item Higher audit costs
		\item Slower reporting
	\end{enumerate}
\end{frame}
\begin{frame}{D1-Q525: Answer}
	\begin{block}{Correct Answer: B}
		Independence is foundational for third-line assurance. Reporting to a management function compromises objectivity and undermines trust in findings.
	\end{block}
\end{frame}

\begin{frame}{D1-Q526: Risk Culture and Incident Reporting}
	\begin{block}{Question}
		The number of reported incidents falls sharply while the risk environment is worsening. What is the MOST likely interpretation?
	\end{block}
	\begin{enumerate}[(A)]
		\item Risk has decreased
		\item A blame culture is suppressing reporting
		\item Controls are effective
		\item The reporting system is too easy to use
	\end{enumerate}
\end{frame}
\begin{frame}{D1-Q526: Answer}
	\begin{block}{Correct Answer: B}
		Falling reports in a worsening environment suggests suppression—classic symptom of a blame culture.
	\end{block}
\end{frame}

\begin{frame}{D1-Q527: Board Risk Education}
	\begin{block}{Question}
		Board members have limited IT risk experience. What should the risk practitioner recommend FIRST?
	\end{block}
	\begin{enumerate}[(A)]
		\item Replace the board
		\item Tailored risk education sessions for board members
		\item Avoid board engagement
		\item Hire an external auditor to educate the board
	\end{enumerate}
\end{frame}
\begin{frame}{D1-Q527: Answer}
	\begin{block}{Correct Answer: B}
		Tailored education builds board capability to oversee risk effectively—essential for governance.
	\end{block}
\end{frame}

\begin{frame}{D1-Q528: Risk Governance in Agile}
	\begin{block}{Question}
		How should risk governance adapt in agile organizations?
	\end{block}
	\begin{enumerate}[(A)]
		\item Abandon risk governance
		\item Embed risk decision-making in agile teams with clear thresholds and escalation
		\item Centralize all risk decisions in a single committee
		\item Delay risk assessment until release
	\end{enumerate}
\end{frame}
\begin{frame}{D1-Q528: Answer}
	\begin{block}{Correct Answer: B}
		Governance must be embedded so teams make risk-aware decisions within thresholds, escalating beyond them.
	\end{block}
\end{frame}

\begin{frame}{D1-Q529: Board IT Subcommittee}
	\begin{block}{Question}
		Which BEST describes the role of a board-level IT risk subcommittee?
	\end{block}
	\begin{enumerate}[(A)]
		\item Execute controls
		\item Provide focused oversight of IT-related risk to the full board
		\item Perform vulnerability scans
		\item Manage the IT budget
	\end{enumerate}
\end{frame}
\begin{frame}{D1-Q529: Answer}
	\begin{block}{Correct Answer: B}
		A subcommittee provides focused oversight of IT-related risk, reporting to the full board with sufficient expertise and attention.
	\end{block}
\end{frame}

\begin{frame}{D1-Q530: Risk Appetite Delegation}
	\begin{block}{Question}
		When delegating risk appetite authority to business units, what is MOST important?
	\end{block}
	\begin{enumerate}[(A)]
		\item That units have documented thresholds and reporting obligations
		\item That units never accept any risk
		\item That authority is undocumented
		\item That units are selected for low turnover
	\end{enumerate}
\end{frame}
\begin{frame}{D1-Q530: Answer}
	\begin{block}{Correct Answer: A}
		Delegated authority requires clear thresholds and reporting so the board can monitor aggregate exposure.
	\end{block}
\end{frame}

\begin{frame}{D1-Q531: Risk Ownership in Shared Services}
	\begin{block}{Question}
		In a shared-services model, who owns the risk of a shared application?
	\end{block}
	\begin{enumerate}[(A)]
		\item The service provider
		\item The business unit that consumes the service
		\item A designated business owner with delegated accountability
		\item No one
	\end{enumerate}
\end{frame}
\begin{frame}{D1-Q531: Answer}
	\begin{block}{Correct Answer: C}
		Shared services still require a designated business owner accountable for the risk and its response.
	\end{block}
\end{frame}

\begin{frame}{D1-Q532: Risk Culture and Leadership}
	\begin{block}{Question}
		Which leadership behavior MOST strongly supports a healthy risk culture?
	\end{block}
	\begin{enumerate}[(A)]
		\item Punishing reporters of issues
		\item Visible modeling of risk-aware decisions and open challenge
		\item Delegating all risk decisions to auditors
		\item Avoiding risk discussions
	\end{enumerate}
\end{frame}
\begin{frame}{D1-Q532: Answer}
	\begin{block}{Correct Answer: B}
		Leaders who model risk awareness and welcome challenge set the tone for a healthy risk culture.
	\end{block}
\end{frame}

\begin{frame}{D1-Q533: Risk Governance Metrics}
	\begin{block}{Question}
		Which governance metric BEST indicates that risk owners are engaged?
	\end{block}
	\begin{enumerate}[(A)]
		\item Number of risk reports
		\item Percentage of risk treatment plans on schedule
		\item Number of software tools
		\item Number of audits
	\end{enumerate}
\end{frame}
\begin{frame}{D1-Q533: Answer}
	\begin{block}{Correct Answer: B}
		On-time treatment plan execution shows owners are actively managing their risks—not just producing paperwork.
	\end{block}
\end{frame}

\begin{frame}{D1-Q534: Risk Governance and Business Changes}
	\begin{block}{Question}
		A major acquisition is planned. What is the MOST important governance action FIRST?
	\end{block}
	\begin{enumerate}[(A)]
		\item Harmonize IT systems
		\item Conduct pre-acquisition risk and governance due diligence
		\item Reduce headcount
		\item Standardize branding
	\end{enumerate}
\end{frame}
\begin{frame}{D1-Q534: Answer}
	\begin{block}{Correct Answer: B}
		Due diligence on the target's risk and governance maturity must come before integration decisions.
	\end{block}
\end{frame}

\begin{frame}{D1-Q535: Risk Reporting and Aggregation}
	\begin{block}{Question}
		Business-unit reports omit cross-unit dependencies. What is the GREATEST consequence?
	\end{block}
	\begin{enumerate}[(A)]
		\item Reports are shorter
		\item Aggregated risk underestimates exposure and may exceed appetite
		\item Vendors are affected
		\item Auditors are inconvenienced
	\end{enumerate}
\end{frame}
\begin{frame}{D1-Q535: Answer}
	\begin{block}{Correct Answer: B}
		Omitting dependencies hides correlated risk and can lead to underestimated exposure at the enterprise level.
	\end{block}
\end{frame}

\begin{frame}{D1-Q536: Governance in Outsourcing}
	\begin{block}{Question}
		What should be established FIRST when outsourcing a critical function?
	\end{block}
	\begin{enumerate}[(A)]
		\item Contract pricing
		\item Governance framework with accountability, reporting, and escalation
		\item Vendor branding
		\item SLA penalties
	\end{enumerate}
\end{frame}
\begin{frame}{D1-Q536: Answer}
	\begin{block}{Correct Answer: B}
		Governance (accountability, reporting, escalation) must be established to oversee outsourced risk—pricing and SLAs follow.
	\end{block}
\end{frame}

\begin{frame}{D1-Q537: Risk Governance and Crisis}
	\begin{block}{Question}
		During a crisis, what is the MOST important governance action?
	\end{block}
	\begin{enumerate}[(A)]
		\item Assign blame
		\item Activate the crisis governance structure with clear decision rights
		\item Wait for auditors
		\item Suppress communication
	\end{enumerate}
\end{frame}
\begin{frame}{D1-Q537: Answer}
	\begin{block}{Correct Answer: B}
		Crisis response requires clear decision rights and communication structures—pre-established for rapid, coordinated action.
	\end{block}
\end{frame}

\begin{frame}{D1-Q538: Board Engagement in Risk}
	\begin{block}{Question}
		What is the BEST indicator that the board is meaningfully engaged in risk oversight?
	\end{block}
	\begin{enumerate}[(A)]
		\item Board meetings occur monthly
		\item Board challenges risk assumptions and integrates risk into strategy decisions
		\item Board receives extensive risk reports
		\item Board has a risk committee
	\end{enumerate}
\end{frame}
\begin{frame}{D1-Q538: Answer}
	\begin{block}{Correct Answer: B}
		Meaningful engagement means challenge, integration of risk into strategy, and informed decisions—not just meetings and reports.
	\end{block}
\end{frame}

\begin{frame}{D1-Q539: Risk Appetite Breach Escalation}
	\begin{block}{Question}
		A critical KRI breaches its threshold. What should happen FIRST?
	\end{block}
	\begin{enumerate}[(A)]
		\item Update the risk register months later
		\item Notify the risk owner for assessment and action
		\item Ignore until the annual review
		\item Conduct a full external audit
	\end{enumerate}
\end{frame}
\begin{frame}{D1-Q539: Answer}
	\begin{block}{Correct Answer: B}
		Immediate notification of the risk owner is required so the breach is assessed and appropriate action initiated.
	\end{block}
\end{frame}

\begin{frame}{D1-Q540: Governance and Corporate Values}
	\begin{block}{Question}
		How should corporate values inform risk governance?
	\end{block}
	\begin{enumerate}[(A)]
		\item They are unrelated
		\item They set expectations for acceptable risk-taking behavior and shape risk appetite
		\item They replace policies
		\item They only matter to auditors
	\end{enumerate}
\end{frame}
\begin{frame}{D1-Q540: Answer}
	\begin{block}{Correct Answer: B}
		Corporate values shape acceptable risk-taking behavior and inform risk appetite, creating coherence between strategy and behavior.
	\end{block}
\end{frame}

% ============================================================
% DOMAIN 2 --- RISK ASSESSMENT: SCENARIO QUESTIONS (Q541–Q565)
% ============================================================
\section{Domain 2: Risk Assessment Scenarios}

\begin{frame}{D2-Q541: New Cloud Migration Assessment}
	\begin{block}{Question}
		Before migrating a critical application to a cloud provider, what must be assessed FIRST?
	\end{block}
	\begin{enumerate}[(A)]
		\item The provider's marketing material
		\item Data classification, regulatory obligations, and provider assurance
		\item The provider's stock price
		\item The number of provider employees
	\end{enumerate}
\end{frame}
\begin{frame}{D2-Q541: Answer}
	\begin{block}{Correct Answer: B}
		Cloud assessment must first consider data classification, regulatory obligations, and provider assurance evidence.
	\end{block}
\end{frame}

\begin{frame}{D2-Q542: Merger Risk Assessment}
	\begin{block}{Question}
		During a merger, which risk is MOST commonly overlooked?
	\end{block}
	\begin{enumerate}[(A)]
		\item Technology incompatibility
		\item Cultural and governance misalignment
		\item Branding conflicts
		\item Office locations
	\end{enumerate}
\end{frame}
\begin{frame}{D2-Q542: Answer}
	\begin{block}{Correct Answer: B}
		Cultural and governance misalignment is consistently the most overlooked integration risk.
	\end{block}
\end{frame}

\begin{frame}{D2-Q543: Risk Assessment for AI}
	\begin{block}{Question}
		What unique risk must be assessed FIRST when deploying AI?
	\end{block}
	\begin{enumerate}[(A)]
		\item Vendor stability
		\item Model transparency, bias, and explainability
		\item Hardware cost
		\item Office location of developers
	\end{enumerate}
\end{frame}
\begin{frame}{D2-Q543: Answer}
	\begin{block}{Correct Answer: B}
		AI-specific risk (bias, explainability, model behavior) must be assessed first, given the unique harms AI can cause.
	\end{block}
\end{frame}

\begin{frame}{D2-Q544: Insider Threat Assessment}
	\begin{block}{Question}
		What is the MOST important factor in assessing insider threat risk?
	\end{block}
	\begin{enumerate}[(A)]
		\item Number of employees
		\item Access privileges combined with behavioral indicators
		\item Office layout
		\item Number of audits
	\end{enumerate}
\end{frame}
\begin{frame}{D2-Q544: Answer}
	\begin{block}{Correct Answer: B}
		Insider risk is driven by access privilege and behavioral signals—the combination is the key assessment factor.
	\end{block}
\end{frame}
	
	
	\begin{frame}{D2-Q545: Risk Assessment for IoT}
		\begin{block}{Question}
			What is the MOST critical assessment focus for IoT devices?
		\end{block}
		\begin{enumerate}[(A)]
			\item Marketing of the IoT vendor
			\item Default credentials, update capability, and network exposure
			\item Number of devices
			\item Power consumption
		\end{enumerate}
	\end{frame}
	\begin{frame}{D2-Q545: Answer}
		\begin{block}{Correct Answer: B}
			IoT default credentials, lack of update mechanisms, and insecure exposure are the dominant sources of risk.
		\end{block}
	\end{frame}
	
	\begin{frame}{D2-Q546: Risk Assessment for OT}
		\begin{block}{Question}
			For operational technology, which risk is MOST unique?
		\end{block}
		\begin{enumerate}[(A)]
			\item Password length
			\item Physical safety consequences of a cyber event
			\item Office ergonomics
			\item Vendor branding
		\end{enumerate}
	\end{frame}
	\begin{frame}{D2-Q546: Answer}
		\begin{block}{Correct Answer: B}
			OT risk can produce physical safety and environmental consequences—unlike traditional IT risk.
		\end{block}
	\end{frame}
	
	\begin{frame}{D2-Q547: Risk Scenario for Ransomware}
		\begin{block}{Question}
			Which is the BEST-formed ransomware risk scenario?
		\end{block}
		\begin{enumerate}[(A)]
			\item "Ransomware might happen"
			\item "Phishing leads to ransomware that encrypts file shares, halting operations for 48 hours"
			\item "Ransomware is bad"
			\item "We need better security"
		\end{enumerate}
	\end{frame}
	\begin{frame}{D2-Q547: Answer}
		\begin{block}{Correct Answer: B}
			A well-formed scenario identifies cause (phishing), event (ransomware), and consequence (operations halted).
		\end{block}
	\end{frame}
	
	\begin{frame}{D2-Q548: Risk Assessment for Supply Chain}
		\begin{block}{Question}
			Which consideration is MOST important in supply chain risk assessment?
		\end{block}
		\begin{enumerate}[(A)]
			\item Number of suppliers
			\item Concentration, single points of failure, and contractual flow-down
			\item Office locations
			\item Vendor branding
		\end{enumerate}
	\end{frame}
	\begin{frame}{D2-Q548: Answer}
		\begin{block}{Correct Answer: B}
			Concentration and contractual flow-down of security requirements determine supply chain resilience.
		\end{block}
	\end{frame}
	
	\begin{frame}{D2-Q549: Risk Assessment for Mergers}
		\begin{block}{Question}
			Which artifact is MOST valuable during MA risk assessment?
		\end{block}
		\begin{enumerate}[(A)]
			\item Marketing deck
			\item Target's risk register, control assessments, and incident history
			\item Brand style guide
			\item Office lease agreements
		\end{enumerate}
	\end{frame}
	\begin{frame}{D2-Q549: Answer}
		\begin{block}{Correct Answer: B}
			The target's risk register, control assessments, and incident history reveal its risk posture and inherited exposures.
		\end{block}
	\end{frame}
	
	\begin{frame}{D2-Q550: Third-Party Assessment Depth}
		\begin{block}{Question}
			Which factor MOST determines depth of a third-party assessment?
		\end{block}
		\begin{enumerate}[(A)]
			\item Vendor size
			\item Criticality of the service and sensitivity of data involved
			\item Vendor location
			\item Length of the contract
		\end{enumerate}
	\end{frame}
	\begin{frame}{D2-Q550: Answer}
		\begin{block}{Correct Answer: B}
			Assessment depth should scale with criticality and data sensitivity—not contract size or vendor location.
		\end{block}
	\end{frame}
	
	\begin{frame}{D2-Q551: Risk Assessment and Business Impact}
		\begin{block}{Question}
			Which is the MOST important output of a BIA for risk assessment?
		\end{block}
		\begin{enumerate}[(A)]
			\item Vendor list
			\item Critical processes, RTOs, RPOs, and impact over time
			\item User survey results
			\item Technology inventory
		\end{enumerate}
	\end{frame}
	\begin{frame}{D2-Q551: Answer}
		\begin{block}{Correct Answer: B}
			BIA outputs critical processes, recovery targets, and impact over time—essential inputs for prioritizing risk.
		\end{block}
	\end{frame}
	
	\begin{frame}{D2-Q552: Risk Register for AI Systems}
		\begin{block}{Question}
			What additional risk register content is needed for AI systems?
		\end{block}
		\begin{enumerate}[(A)]
			\item Vendor name only
			\item Model purpose, training data, and explainability information
			\item Code lines of code
			\item Office location
		\end{enumerate}
	\end{frame}
	\begin{frame}{D2-Q552: Answer}
		\begin{block}{Correct Answer: B}
			AI risk register entries require model purpose, training data provenance, and explainability records.
		\end{block}
	\end{frame}
	
	\begin{frame}{D2-Q553: Assessment of New Tech}
		\begin{block}{Question}
			When introducing a novel technology, what should be done FIRST?
		\end{block}
		\begin{enumerate}[(A)]
			\item Pilot deployment to all users
			\item Threat modeling and risk assessment
			\item Purchase licenses
			\item Announce to media
		\end{enumerate}
	\end{frame}
	\begin{frame}{D2-Q553: Answer}
		\begin{block}{Correct Answer: B}
			Threat modeling and risk assessment of novel technology should precede any deployment.
		\end{block}
	\end{frame}
	
	\begin{frame}{D2-Q554: Risk Assessment Iteration}
		\begin{block}{Question}
			How often should the risk register be fully reassessed?
		\end{block}
		\begin{enumerate}[(A)]
			\item Never
			\item At a cadence matched to the rate of change in the environment and criticality
			\item Only after incidents
			\item Only when auditors ask
		\end{enumerate}
	\end{frame}
	\begin{frame}{D2-Q554: Answer}
		\begin{block}{Correct Answer: B}
			Reassessment cadence should be proportional to the pace of change and the criticality of what is assessed.
		\end{block}
	\end{frame}
	
	\begin{frame}{D2-Q555: Assessment of Business Changes}
		\begin{block}{Question}
			What should trigger an immediate ad hoc risk assessment?
		\end{block}
		\begin{enumerate}[(A)]
			\item A new hire
			\item A significant business change (e.g., new product, geography, or technology)
			\item A new auditor
			\item A vendor rebrand
		\end{enumerate}
	\end{frame}
	\begin{frame}{D2-Q555: Answer}
		\begin{block}{Correct Answer: B}
			Significant business changes alter risk exposure and require prompt reassessment.
		\end{block}
	\end{frame}
	
	\begin{frame}{D2-Q556: Risk Assessment and Insurance}
		\begin{block}{Question}
			How does risk assessment inform insurance decisions?
		\end{block}
		\begin{enumerate}[(A)]
			\item It has no role
			\item It quantifies exposure so insurance limits and deductibles can be set
			\item It only matters to vendors
			\item It should be ignored
		\end{enumerate}
	\end{frame}
	\begin{frame}{D2-Q556: Answer}
		\begin{block}{Correct Answer: B}
			Risk assessment quantifies exposure, enabling appropriate insurance limits, deductibles, and coverage selection.
		\end{block}
	\end{frame}
	
	\begin{frame}{D2-Q557: Risk Assessment and Business Continuity}
		\begin{block}{Question}
			Which is the PRIMARY link between risk assessment and business continuity?
		\end{block}
		\begin{enumerate}[(A)]
			\item Vendor selection
			\item Risk assessment identifies threats and impacts; BIA quantifies recovery needs; together they drive BC planning
			\item Neither is related
			\item Both are optional
		\end{enumerate}
	\end{frame}
	\begin{frame}{D2-Q557: Answer}
		\begin{block}{Correct Answer: B}
			Risk assessment identifies threats and impacts; BIA quantifies recovery needs. Combined, they inform BC strategy and testing.
		\end{block}
	\end{frame}
	
	\begin{frame}{D2-Q558: Risk Assessment for Cloud Misconfiguration}
		\begin{block}{Question}
			What is the leading cloud risk category?
		\end{block}
		\begin{enumerate}[(A)]
			\item Provider bankruptcies
			\item Customer misconfiguration of cloud resources
			\item Data center weather events
			\item Provider stock volatility
		\end{enumerate}
	\end{frame}
	\begin{frame}{D2-Q558: Answer}
		\begin{block}{Correct Answer: B}
			Customer misconfiguration (public buckets, weak IAM) is consistently the leading cloud risk.
		\end{block}
	\end{frame}
	
	\begin{frame}{D2-Q559: Risk Assessment and Maturity}
		\begin{block}{Question}
			How does risk management maturity affect risk assessment quality?
		\end{block}
		\begin{enumerate}[(A)]
			\item Maturity has no effect
			\item Higher maturity yields more consistent, documented, and repeatable assessments
			\item Lower maturity yields better assessments
			\item Maturity equals zero risk
		\end{enumerate}
	\end{frame}
	\begin{frame}{D2-Q559: Answer}
		\begin{block}{Correct Answer: B}
			Maturity correlates with consistent, documented, and repeatable assessment practices.
		\end{block}
	\end{frame}
	
	\begin{frame}{D2-Q560: Assessment of Third-Party Concentration}
		\begin{block}{Question}
			Why is vendor concentration a key risk assessment concern?
		\end{block}
		\begin{enumerate}[(A)]
			\item Vendor branding
			\item A single provider disruption can affect multiple critical services simultaneously
			\item Vendor office location
			\item Vendor marketing
		\end{enumerate}
	\end{frame}
	\begin{frame}{D2-Q560: Answer}
		\begin{block}{Correct Answer: B}
			Concentration creates correlated exposure—one provider failure can disrupt multiple services.
		\end{block}
	\end{frame}
	
	\begin{frame}{D2-Q561: Risk Assessment and Regulatory Change}
		\begin{block}{Question}
			A new privacy regulation passes. What is the FIRST risk assessment action?
		\end{block}
		\begin{enumerate}[(A)]
			\item Update policies immediately
			\item Assess impact on business processes, data flows, and existing controls
			\item Ignore until enforcement
			\item Announce compliance publicly
		\end{enumerate}
	\end{frame}
	\begin{frame}{D2-Q561: Answer}
		\begin{block}{Correct Answer: B}
			Impact assessment comes first—policies and controls should follow an understanding of the regulation's effects.
		\end{block}
	\end{frame}
	
	\begin{frame}{D2-Q562: Risk Assessment and Scenario Planning}
		\begin{block}{Question}
			What is the PRIMARY benefit of scenario planning in risk assessment?
		\end{block}
		\begin{enumerate}[(A)]
			\item It replaces the risk register
			\item It surfaces non-obvious exposures and dependencies
			\item It eliminates risk
			\item It satisfies regulators
		\end{enumerate}
	\end{frame}
	\begin{frame}{D2-Q562: Answer}
		\begin{block}{Correct Answer: B}
			Scenario planning surfaces exposures and dependencies that checklists miss.
		\end{block}
	\end{frame}
	
	\begin{frame}{D2-Q563: Risk Assessment and Insider Risk}
		\begin{block}{Question}
			Which is the MOST effective control for insider risk after assessment?
		\end{block}
		\begin{enumerate}[(A)]
			\item Firewalls
			\item Least privilege combined with monitoring and behavioral analytics
			\item Antivirus
			\item Strong passwords only
		\end{enumerate}
	\end{frame}
	\begin{frame}{D2-Q563: Answer}
		\begin{block}{Correct Answer: B}
			Insider risk is best managed through least privilege plus monitoring and behavioral analytics.
		\end{block}
	\end{frame}
	
	\begin{frame}{D2-Q564: Risk Assessment and Change}
		\begin{block}{Question}
			Why must risk assessment be repeated after significant change?
		\end{block}
		\begin{enumerate}[(A)]
			\item To keep auditors busy
			\item Because change alters exposure, dependencies, and control effectiveness
			\item Because regulators require it
			\item Because vendors ask
		\end{enumerate}
	\end{frame}
	\begin{frame}{D2-Q564: Answer}
		\begin{block}{Correct Answer: B}
			Change alters exposure and control effectiveness—reassessment is required to maintain accurate risk understanding.
		\end{block}
	\end{frame}
	
	\begin{frame}{D2-Q565: Risk Assessment Output for Board}
		\begin{block}{Question}
			Which is the MOST valuable risk assessment output for the board?
		\end{block}
		\begin{enumerate}[(A)]
			\item Raw vulnerability lists
			\item Aggregated risk against appetite with clear recommendations
			\item Detailed network diagrams
			\item All technical findings
		\end{enumerate}
	\end{frame}
	\begin{frame}{D2-Q565: Answer}
		\begin{block}{Correct Answer: B}
			The board needs aggregated risk vs. appetite with clear recommendations—technical detail obscures oversight.
		\end{block}
	\end{frame}
	
	% ============================================================
	% DOMAIN 3 --- RISK RESPONSE AND REPORTING SCENARIOS (Q566–Q595)
	% ============================================================
	\section{Domain 3: Risk Response Scenarios}
	
	\begin{frame}{D3-Q566: Cloud Incident Response}
		\begin{block}{Question}
			What is MOST important to establish with a cloud provider before an incident?
		\end{block}
		\begin{enumerate}[(A)]
			\item Marketing contacts
			\item Incident notification and forensic cooperation requirements
			\item Stock price disclosure
			\item Office locations
		\end{enumerate}
	\end{frame}
	\begin{frame}{D3-Q566: Answer}
		\begin{block}{Correct Answer: B}
			Incident notification SLAs and forensic cooperation must be contractually established in advance.
		\end{block}
	\end{frame}
	
	\begin{frame}{D3-Q567: Ransomware Response}
		\begin{block}{Question}
			What is the MOST important FIRST response action to ransomware?
		\end{block}
		\begin{enumerate}[(A)]
			\item Pay the ransom
			\item Isolate affected systems and notify the incident response team
			\item Communicate to press
			\item Blame employees
		\end{enumerate}
	\end{frame}
	\begin{frame}{D3-Q567: Answer}
		\begin{block}{Correct Answer: B}
			Containment (isolation) and IR activation precede all other actions—including ransom decisions.
		\end{block}
	\end{frame}
	
	\begin{frame}{D3-Q568: Third-Party Breach Response}
		\begin{block}{Question}
			A vendor reports a breach. What should the enterprise do FIRST?
		\end{block}
		\begin{enumerate}[(A)]
			\item Terminate the contract
			\item Activate incident coordination, assess data exposure, and notify regulators as required
			\item Ignore
			\item Announce to press
		\end{enumerate}
	\end{frame}
	\begin{frame}{D3-Q568: Answer}
		\begin{block}{Correct Answer: B}
			Coordination, data-exposure assessment, and required notifications are the correct immediate actions.
		\end{block}
	\end{frame}
	
	\begin{frame}{D3-Q569: Insider Incident Response}
		\begin{block}{Question}
			For a suspected insider data theft, what is MOST critical FIRST?
		\end{block}
		\begin{enumerate}[(A)]
			\item Confront the employee
			\item Preserve evidence and involve legal/HR, avoiding tipping off the suspect
			\item Ignore
			\item Announce to press
		\end{enumerate}
	\end{frame}
	\begin{frame}{D3-Q569: Answer}
		\begin{block}{Correct Answer: B}
			Preserving evidence and coordinating with legal/HR—without tipping off—is essential to investigation success.
		\end{block}
	\end{frame}
	
	\begin{frame}{D3-Q570: Business Continuity Activation}
		\begin{block}{Question}
			When should a BCP be activated?
		\end{block}
		\begin{enumerate}[(A)]
			\item Whenever any incident occurs
			\item When a disruption exceeds defined activation criteria and thresholds
			\item Only after regulators demand it
			\item Only annually
		\end{enumerate}
	\end{frame}
	\begin{frame}{D3-Q570: Answer}
		\begin{block}{Correct Answer: B}
			Activation should follow predefined criteria and thresholds—not every incident warrants full continuity activation.
		\end{block}
	\end{frame}
	
	\begin{frame}{D3-Q571: Communications During Crisis}
		\begin{block}{Question}
			What is MOST important about crisis communications?
		\end{block}
		\begin{enumerate}[(A)]
			\item Speeding to press without facts
			\item Preapproved channels, spokespeople, and templates with factual updates
			\item Silence
			\item Blaming vendors
		\end{enumerate}
	\end{frame}
	\begin{frame}{D3-Q571: Answer}
		\begin{block}{Correct Answer: B}
			Preapproved channels, spokespeople, and templates ensure clear, factual, and controlled communications.
		\end{block}
	\end{frame}
	
	\begin{frame}{D3-Q572: Risk Treatment for Legacy Systems}
		\begin{block}{Question}
			A critical legacy system cannot be patched. What is the BEST response?
		\end{block}
		\begin{enumerate}[(A)]
			\item Ignore
			\item Compensating controls (segmentation, monitoring) with documented acceptance
			\item Turn off the system
			\item Blame the vendor
		\end{enumerate}
	\end{frame}
	\begin{frame}{D3-Q572: Answer}
		\begin{block}{Correct Answer: B}
			Compensating controls plus documented acceptance are appropriate when patching is infeasible.
		\end{block}
	\end{frame}
	
	\begin{frame}{D3-Q573: Risk Response for Remote Work}
		\begin{block}{Question}
			What is the MOST important remote-work risk response?
		\end{block}
		\begin{enumerate}[(A)]
			\item VPN only
			\item Layered defenses: device management, zero trust, MFA, and monitoring
			\item Password changes only
			\item Buying new laptops
		\end{enumerate}
	\end{frame}
	\begin{frame}{D3-Q573: Answer}
		\begin{block}{Correct Answer: B}
			Remote work requires layered defenses—device management, zero trust, MFA, and monitoring—not a single control.
		\end{block}
	\end{frame}
	
	\begin{frame}{D3-Q574: Risk Response for Shadow IT}
		\begin{block}{Question}
			What is the BEST response to Shadow IT?
		\end{block}
		\begin{enumerate}[(A)]
			\item Punish users
			\item Discover, assess, and either secure or provide approved alternatives
			\item Ignore
			\item Block everything
		\end{enumerate}
	\end{frame}
	\begin{frame}{D3-Q574: Answer}
		\begin{block}{Correct Answer: B}
			Discovery plus assessment and approved alternatives addresses the underlying need—purely blocking often fails.
		\end{block}
	\end{frame}
	
	\begin{frame}{D3-Q575: Response to Vulnerability Exploit}
		\begin{block}{Question}
			A critical vulnerability is being actively exploited. What should happen FIRST?
		\end{block}
		\begin{enumerate}[(A)]
			\item Wait for the next patch cycle
			\item Apply emergency mitigation (compensating controls or emergency patch)
			\item Ignore
			\item Announce to press
		\end{enumerate}
	\end{frame}
	\begin{frame}{D3-Q575: Answer}
		\begin{block}{Correct Answer: B}
			Active exploitation requires emergency mitigation—not waiting for the next patch cycle.
		\end{block}
	\end{frame}
	
	\begin{frame}{D3-Q576: Risk Response for MA}
		\begin{block}{Question}
			What is the MOST important post-merger risk response action?
		\end{block}
		\begin{enumerate}[(A)]
			\item Rebrand quickly
			\item Integrate governance, controls, and risk registers
			\item Lay off employees
			\item Standardize office furniture
		\end{enumerate}
	\end{frame}
	\begin{frame}{D3-Q576: Answer}
		\begin{block}{Correct Answer: B}
			Integrating governance and controls is essential to realize the value of the merger and manage inherited risk.
		\end{block}
	\end{frame}
	
	\begin{frame}{D3-Q577: Risk Response for AI}
		\begin{block}{Question}
			What is the MOST important AI risk response control?
		\end{block}
		\begin{enumerate}[(A)]
			\item Marketing
			\item Human oversight, monitoring, and bias testing
			\item Faster hardware
			\item Office location
		\end{enumerate}
	\end{frame}
	\begin{frame}{D3-Q577: Answer}
		\begin{block}{Correct Answer: B}
			Human oversight, ongoing monitoring, and bias testing are essential AI risk controls.
		\end{block}
	\end{frame}
	
	\begin{frame}{D3-Q578: Risk Response for Data Loss}
		\begin{block}{Question}
			What is the MOST effective layered defense for data loss?
		\end{block}
		\begin{enumerate}[(A)]
			\item Firewall alone
			\item DLP, encryption, access controls, and monitoring together
			\item Antivirus alone
			\item Password changes
		\end{enumerate}
	\end{frame}
	\begin{frame}{D3-Q578: Answer}
		\begin{block}{Correct Answer: B}
			DLP, encryption, access controls, and monitoring together form an effective layered defense.
		\end{block}
	\end{frame}
	
	\begin{frame}{D3-Q579: Risk Response and Vendor Exit}
		\begin{block}{Question}
			What is MOST important in a vendor exit strategy?
		\end{block}
		\begin{enumerate}[(A)]
			\item Branding
			\item Data return/destruction, transition support, and continuity provisions
			\item Marketing
			\item Office location
		\end{enumerate}
	\end{frame}
	\begin{frame}{D3-Q579: Answer}
		\begin{block}{Correct Answer: B}
			Exit strategies must specify data return/destruction, transition support, and continuity provisions.
		\end{block}
	\end{frame}
	
	\begin{frame}{D3-Q580: Risk Response and Regulatory Investigation}
		\begin{block}{Question}
			Regulators begin an investigation. What should happen FIRST?
		\end{block}
		\begin{enumerate}[(A)]
			\item Speak publicly
			\item Engage legal counsel and preserve evidence
			\item Delete records
			\item Blame the vendor
		\end{enumerate}
	\end{frame}
	\begin{frame}{D3-Q580: Answer}
		\begin{block}{Correct Answer: B}
			Legal counsel engagement and evidence preservation are the correct first steps in a regulatory investigation.
		\end{block}
	\end{frame}
	
	\begin{frame}{D3-Q581: Risk Response for Misconfiguration}
		\begin{block}{Question}
			The leading cause of cloud breaches is misconfiguration. What is the BEST response?
		\end{block}
		\begin{enumerate}[(A)]
			\item Ignore
			\item Baseline configurations, automated compliance scanning, and change control
			\item Manual review annually
			\item Password changes
		\end{enumerate}
	\end{frame}
	\begin{frame}{D3-Q581: Answer}
		\begin{block}{Correct Answer: B}
			Baselines plus automated scanning and change control address misconfiguration at scale.
		\end{block}
	\end{frame}
	
	\begin{frame}{D3-Q582: Risk Response for Third Party}
		\begin{block}{Question}
			What is the MOST effective control for third-party risk?
		\end{block}
		\begin{enumerate}[(A)]
			\item Vendor marketing
			\item Contractual requirements plus independent assurance and monitoring
			\item Vendor size
			\item Vendor location
		\end{enumerate}
	\end{frame}
	\begin{frame}{D3-Q582: Answer}
		\begin{block}{Correct Answer: B}
			Contractual controls plus independent assurance and ongoing monitoring constitute the strongest third-party control.
		\end{block}
	\end{frame}
	
	\begin{frame}{D3-Q583: Response to Breach Notification}
		\begin{block}{Question}
			After a breach notification, what is the MOST important action?
		\end{block}
		\begin{enumerate}[(A)]
			\item Hide it
			\item Assess scope, notify as required, and implement lessons learned
			\item Blame employees
			\item Delay
		\end{enumerate}
	\end{frame}
	\begin{frame}{D3-Q583: Answer}
		\begin{block}{Correct Answer: B}
			Scope assessment, required notifications, and lessons learned are the correct post-breach actions.
		\end{block}
	\end{frame}
	
	\begin{frame}{D3-Q584: Risk Response and Business Change}
		\begin{block}{Question}
			After a major business change, what should be revisited FIRST?
		\end{block}
		\begin{enumerate}[(A)]
			\item Brand guidelines
			\item Risk assessment and treatment plan alignment
			\item Office layout
			\item Marketing
		\end{enumerate}
	\end{frame}
	\begin{frame}{D3-Q584: Answer}
		\begin{block}{Correct Answer: B}
			Risk assessments and treatment plans must be revisited and realigned after significant business change.
		\end{block}
	\end{frame}
	
	\begin{frame}{D3-Q585: Risk Response and Residual Risk}
		\begin{block}{Question}
			Residual risk remains above tolerance after treatment. What should happen?
		\end{block}
		\begin{enumerate}[(A)]
			\item Ignore
			\item Escalate to the risk owner for further treatment or formal acceptance
			\item Adjust the assessment
			\item Blame the vendor
		\end{enumerate}
	\end{frame}
	\begin{frame}{D3-Q585: Answer}
		\begin{block}{Correct Answer: B}
			Residual risk above tolerance must be escalated for further treatment or formal acceptance by the appropriate owner.
		\end{block}
	\end{frame}
	
	\begin{frame}{D3-Q586: Risk Response and Insurance Renewal}
		\begin{block}{Question}
			Before renewing cyber insurance, what is MOST important?
		\end{block}
		\begin{enumerate}[(A)]
			\item Compare premium only
			\item Confirm coverage aligns to current risk exposure and control posture
			\item Negotiate faster payment
			\item Ignore exclusions
		\end{enumerate}
	\end{frame}
	\begin{frame}{D3-Q586: Answer}
		\begin{block}{Correct Answer: B}
			Coverage should be reviewed against current exposure and controls—premium alone is an insufficient decision basis.
		\end{block}
	\end{frame}
	
	\begin{frame}{D3-Q587: Risk Response and Crisis}
		\begin{block}{Question}
			During a crisis, what is MOST important about decision-making?
		\end{block}
		\begin{enumerate}[(A)]
			\item Slow and hierarchical
			\item Clear delegated authority and pre-approved playbooks
			\item Consensus at every step
			\item Avoid communication
		\end{enumerate}
	\end{frame}
	\begin{frame}{D3-Q587: Answer}
		\begin{block}{Correct Answer: B}
			Crisis decisions require pre-delegated authority and playbooks to act quickly and coherently.
		\end{block}
	\end{frame}
	
	\begin{frame}{D3-Q588: Risk Response and Lessons Learned}
		\begin{block}{Question}
			After an incident, what is MOST important about lessons learned?
		\end{block}
		\begin{enumerate}[(A)]
			\item Blame individuals
			\item Systematically update policies, controls, and playbooks
			\item File and forget
			\item Announce publicly
		\end{enumerate}
	\end{frame}
	\begin{frame}{D3-Q588: Answer}
		\begin{block}{Correct Answer: B}
			Lessons learned must translate into updates to policy, control, and playbooks—not blame.
		\end{block}
	\end{frame}
	
	\begin{frame}{D3-Q589: Risk Response and Business Continuity Testing}
		\begin{block}{Question}
			Which test is MOST valuable for a newly created BCP?
		\end{block}
		\begin{enumerate}[(A)]
			\item Document review
			\item Tabletop exercise followed by targeted technical tests
			\item Marketing review
			\item No testing
		\end{enumerate}
	\end{frame}
	\begin{frame}{D3-Q589: Answer}
		\begin{block}{Correct Answer: B}
			Tabletop exercises plus targeted technical tests validate and improve the plan.
		\end{block}
	\end{frame}
	
	\begin{frame}{D3-Q590: Risk Response and Shared Responsibility}
		\begin{block}{Question}
			Under the shared responsibility model, what must be defined FIRST?
		\end{block}
		\begin{enumerate}[(A)]
			\item Pricing
			\item Which controls each party is responsible for
			\item Marketing
			\item Office location
		\end{enumerate}
	\end{frame}
	\begin{frame}{D3-Q590: Answer}
		\begin{block}{Correct Answer: B}
			Explicit definition of which controls belong to which party is essential to prevent gaps and overlaps.
		\end{block}
	\end{frame}
	
	\begin{frame}{D3-Q591: Risk Response and Compliance}
		\begin{block}{Question}
			What is the MOST important compliance risk response control?
		\end{block}
		\begin{enumerate}[(A)]
			\item Marketing
			\item Continuous compliance monitoring against regulatory obligations
			\item Annual audit alone
			\item Password changes
		\end{enumerate}
	\end{frame}
	\begin{frame}{D3-Q591: Answer}
		\begin{block}{Correct Answer: B}
			Continuous compliance monitoring provides real-time assurance; annual audits alone leave gaps.
		\end{block}
	\end{frame}
	
	\begin{frame}{D3-Q592: Risk Response for Supply Chain}
		\begin{block}{Question}
			What is the most important supply chain risk response?
		\end{block}
		\begin{enumerate}[(A)]
			\item Vendor branding
			\item Contractual security requirements plus assurance and diversification
			\item Vendor location
			\item Vendor size
		\end{enumerate}
	\end{frame}
	\begin{frame}{D3-Q592: Answer}
		\begin{block}{Correct Answer: B}
			Contractual controls, independent assurance, and diversification reduce supply chain risk materially.
		\end{block}
	\end{frame}
	
	\begin{frame}{D3-Q593: Risk Response and Third-Party Access}
		\begin{block}{Question}
			What is the most important control for third-party access to systems?
		\end{block}
		\begin{enumerate}[(A)]
			\item Password complexity alone
			\item Least privilege, MFA, monitoring, and timely revocation
			\item Marketing
			\item Office location
		\end{enumerate}
	\end{frame}
	\begin{frame}{D3-Q593: Answer}
		\begin{block}{Correct Answer: B}
			Least privilege, MFA, monitoring, and timely revocation are the essential controls for third-party access.
		\end{block}
	\end{frame}
	
	\begin{frame}{D3-Q594: Risk Response and Regulatory Penalties}
		\begin{block}{Question}
			A regulator proposes a significant fine. What is the BEST response?
		\end{block}
		\begin{enumerate}[(A)]
			\item Ignore
			\item Engage counsel, remediate root causes, and negotiate remediation credit where possible
			\item Blame vendors
			\item Announce to press without counsel
		\end{enumerate}
	\end{frame}
	\begin{frame}{D3-Q594: Answer}
		\begin{block}{Correct Answer: B}
			Counsel engagement, root-cause remediation, and negotiating remediation credit are the correct steps.
		\end{block}
	\end{frame}
	
	\begin{frame}{D3-Q595: Risk Response and Continuous Improvement}
		\begin{block}{Question}
			What is the MOST important continuous improvement activity after a risk event?
		\end{block}
		\begin{enumerate}[(A)]
			\item Blame
			\item Update risk assessments, controls, and playbooks based on outcomes
			\item Forget
			\item Increase premiums
		\end{enumerate}
	\end{frame}
	\begin{frame}{D3-Q595: Answer}
		\begin{block}{Correct Answer: B}
			Systematic updates to assessments, controls, and playbooks embed learning into the risk program.
		\end{block}
	\end{frame}
	
	% ============================================================
	% DOMAIN 4 --- IT AND SECURITY SCENARIOS (Q596–Q615)
	% ============================================================
	\section{Domain 4: IT and Security Scenarios}
	
	\begin{frame}{D4-Q596: Securing Remote Access}
		\begin{block}{Question}
			What is the MOST important control for remote access security?
		\end{block}
		\begin{enumerate}[(A)]
			\item Long passwords
			\item Zero trust with MFA and least privilege
			\item Single VPN
			\item Office location
		\end{enumerate}
	\end{frame}
	\begin{frame}{D4-Q596: Answer}
		\begin{block}{Correct Answer: B}
			Zero trust with MFA and least privilege—not perimeter-only controls—secures remote access.
		\end{block}
	\end{frame}
	
	\begin{frame}{D4-Q597: Protecting Data at Rest}
		\begin{block}{Question}
			What is the most effective control for data at rest?
		\end{block}
		\begin{enumerate}[(A)]
			\item Hashing
			\item Encryption with strong key management
			\item Password changes
			\item Firewall
		\end{enumerate}
	\end{frame}
	\begin{frame}{D4-Q597: Answer}
		\begin{block}{Correct Answer: B}
			Encryption with strong key management protects data at rest; hashing does not provide confidentiality.
		\end{block}
	\end{frame}
	
	\begin{frame}{D4-Q598: Protecting Data in Transit}
		\begin{block}{Question}
			What is the most effective control for data in transit?
		\end{block}
		\begin{enumerate}[(A)]
			\item Hashing
			\item TLS with strong cipher suites
			\item Antivirus
			\item Password changes
		\end{enumerate}
	\end{frame}
	\begin{frame}{D4-Q598: Answer}
		\begin{block}{Correct Answer: B}
			TLS with strong cipher suites is the standard control for protecting data in transit.
		\end{block}
	\end{frame}
	
	\begin{frame}{D4-Q599: Identity Provider Compromise}
		\begin{block}{Question}
			What is the GREATEST risk if an identity provider is compromised?
		\end{block}
		\begin{enumerate}[(A)]
			\item Increased workload
			\item Widespread access compromise across federated systems
			\item Vendor branding
			\item Office location
		\end{enumerate}
	\end{frame}
	\begin{frame}{D4-Q599: Answer}
		\begin{block}{Correct Answer: B}
			A compromised identity provider can enable access to all federated systems—an enterprise-wide single point of failure.
		\end{block}
	\end{frame}
	
	\begin{frame}{D4-Q600: Application Security}
		\begin{block}{Question}
			What is the most effective way to reduce application-level vulnerabilities?
		\end{block}
		\begin{enumerate}[(A)]
			\item Weekly password changes
			\item Secure coding practices, SAST/DAST, and dependency management
			\item Firewall only
			\item Office location
		\end{enumerate}
	\end{frame}
	\begin{frame}{D4-Q600: Answer}
		\begin{block}{Correct Answer: B}
			Secure coding practices, SAST/DAST, and dependency management reduce application vulnerabilities at scale.
		\end{block}
	\end{frame}
	
	\begin{frame}{D4-Q601: Protecting Against Phishing}
		\begin{block}{Question}
			What is the MOST effective technical control against phishing?
		\end{block}
		\begin{enumerate}[(A)]
			\item Longer passwords
			\item DMARC, URL filtering, and MFA on accounts
			\item Firewalls only
			\item Office location
		\end{enumerate}
	\end{frame}
	\begin{frame}{D4-Q601: Answer}
		\begin{block}{Correct Answer: B}
			DMARC, URL filtering, and MFA materially reduce phishing impact and credential compromise.
		\end{block}
	\end{frame}
	
	\begin{frame}{D4-Q602: Data Classification Challenge}
		\begin{block}{Question}
			What is the MAIN challenge in data classification?
		\end{block}
		\begin{enumerate}[(A)]
			\item Too few files
			\item Consistency across systems and lifecycle stages
			\item Encryption
			\item Firewalls
		\end{enumerate}
	\end{frame}
	\begin{frame}{D4-Q602: Answer}
		\begin{block}{Correct Answer: B}
			Maintaining consistent classification across systems and lifecycle stages is the primary challenge.
		\end{block}
	\end{frame}
	
	\begin{frame}{D4-Q603: Third-Party Access Control}
		\begin{block}{Question}
			What is the most important control for third-party access?
		\end{block}
		\begin{enumerate}[(A)]
			\item Long contracts
			\item Time-bound, least-privilege access with monitoring and timely revocation
			\item Marketing
			\item Office location
		\end{enumerate}
	\end{frame}
	\begin{frame}{D4-Q603: Answer}
		\begin{block}{Correct Answer: B}
			Time-bound, least-privilege access with monitoring and revocation is essential for third-party access security.
		\end{block}
	\end{frame}
	
	\begin{frame}{D4-Q604: Cloud Misconfiguration}
		\begin{block}{Question}
			What is the most effective way to reduce cloud misconfiguration risk?
		\end{block}
		\begin{enumerate}[(A)]
			\item Longer passwords
			\item Infrastructure-as-code baselines with automated scanning
			\item Firewalls only
			\item Office location
		\end{enumerate}
	\end{frame}
	\begin{frame}{D4-Q604: Answer}
		\begin{block}{Correct Answer: B}
			IaC baselines with automated scanning reduce misconfiguration at scale more effectively than manual review.
		\end{block}
	\end{frame}
	
	\begin{frame}{D4-Q605: Security Awareness Measurement}
		\begin{block}{Question}
			What is the BEST way to measure security awareness effectiveness?
		\end{block}
		\begin{enumerate}[(A)]
			\item Number of emails sent
			\item Simulation results, reporting rates, and behavioral metrics over time
			\item Posters distributed
			\item Office location
		\end{enumerate}
	\end{frame}
	\begin{frame}{D4-Q605: Answer}
		\begin{block}{Correct Answer: B}
			Simulation results, reporting rates, and behavioral metrics over time are the best evidence of awareness effectiveness.
		\end{block}
	\end{frame}
	
	\begin{frame}{D4-Q606: Data Breach Response}
		\begin{block}{Question}
			What is the FIRST action after discovering a breach?
		\end{block}
		\begin{enumerate}[(A)]
			\item Announce to press
			\item Contain the breach and activate incident response
			\item Blame employees
			\item Delay
		\end{enumerate}
	\end{frame}
	\begin{frame}{D4-Q606: Answer}
		\begin{block}{Correct Answer: B}
			Containment and IR activation come first—disclosure and other steps follow assessment.
		\end{block}
	\end{frame}
	
	\begin{frame}{D4-Q607: Insider Threat Program}
		\begin{block}{Question}
			What is the most important element of an insider threat program?
		\end{block}
		\begin{enumerate}[(A)]
			\item Marketing
			\item Cross-functional detection, investigation, and legal coordination
			\item Firewall alone
			\item Office location
		\end{enumerate}
	\end{frame}
	\begin{frame}{D4-Q607: Answer}
		\begin{block}{Correct Answer: B}
			Insider threat programs require coordinated detection, investigation, and legal/HR involvement.
		\end{block}
	\end{frame}
	
	\begin{frame}{D4-Q608: Patch Management Strategy}
		\begin{block}{Question}
			What is the MOST important factor in patch management strategy?
		\end{block}
		\begin{enumerate}[(A)]
			\item Marketing
			\item Risk-based prioritization with testing and rollback
			\item Fastest patch application
			\item Office location
		\end{enumerate}
	\end{frame}
	\begin{frame}{D4-Q608: Answer}
		\begin{block}{Correct Answer: B}
			Risk-based prioritization with testing and rollback balances speed and stability.
		\end{block}
	\end{frame}
	
	\begin{frame}{D4-Q609: Log Management Program}
		\begin{block}{Question}
			What is the MOST important element of a log management program?
		\end{block}
		\begin{enumerate}[(A)]
			\item Volume of logs
			\item Centralized, protected, and actively reviewed logs per policy
			\item Marketing
			\item Office location
		\end{enumerate}
	\end{frame}
	\begin{frame}{D4-Q609: Answer}
		\begin{block}{Correct Answer: B}
			Logs must be centralized, protected, and actively reviewed—collection alone provides no value.
		\end{block}
	\end{frame}
	
	\begin{frame}{D4-Q610: Security Architecture Review}
		\begin{block}{Question}
			What is the PRIMARY purpose of a security architecture review?
		\end{block}
		\begin{enumerate}[(A)]
			\item Marketing
			\item Identify design gaps and misalignments with business and risk requirements
			\item Office location
			\item Vendor branding
		\end{enumerate}
	\end{frame}
	\begin{frame}{D4-Q610: Answer}
		\begin{block}{Correct Answer: B}
			Security architecture reviews identify design gaps and misalignment with business and risk requirements.
		\end{block}
	\end{frame}
	
	\begin{frame}{D4-Q611: Multi-Cloud Security}
		\begin{block}{Question}
			What is the MOST important multi-cloud security consideration?
		\end{block}
		\begin{enumerate}[(A)]
			\item Vendor branding
			\item Consistent policy, identity, and monitoring across clouds
			\item Office location
			\item Vendor size
		\end{enumerate}
	\end{frame}
	\begin{frame}{D4-Q611: Answer}
		\begin{block}{Correct Answer: B}
			Consistent policy, identity, and monitoring prevent gaps and inconsistencies in multi-cloud environments.
		\end{block}
	\end{frame}
	
	\begin{frame}{D4-Q612: Artificial Intelligence Security}
		\begin{block}{Question}
			What is a key AI-specific security risk?
		\end{block}
		\begin{enumerate}[(A)]
			\item Marketing
			\item Adversarial inputs and data poisoning
			\item Office location
			\item Vendor branding
		\end{enumerate}
	\end{frame}
	\begin{frame}{D4-Q612: Answer}
		\begin{block}{Correct Answer: B}
			AI is uniquely exposed to adversarial inputs and data poisoning—risks absent in traditional systems.
		\end{block}
	\end{frame}
	
	\begin{frame}{D4-Q613: OT Security}
		\begin{block}{Question}
			What is the MOST important OT security consideration?
		\end{block}
		\begin{enumerate}[(A)]
			\item Marketing
			\item Safety and availability with strict change control
			\item Office location
			\item Vendor branding
		\end{enumerate}
	\end{frame}
	\begin{frame}{D4-Q613: Answer}
		\begin{block}{Correct Answer: B}
			OT security prioritizes safety and availability with strict change control—unlike IT where confidentiality is often paramount.
		\end{block}
	\end{frame}
	
	\begin{frame}{D4-Q614: Third-Party SOC Reports}
		\begin{block}{Question}
			Which SOC report provides the STRONGEST assurance of ongoing control effectiveness?
		\end{block}
		\begin{enumerate}[(A)]
			\item SOC 1 Type I
			\item SOC 2 Type II
			\item Marketing
			\item SOC 3
		\end{enumerate}
	\end{frame}
	\begin{frame}{D4-Q614: Answer}
		\begin{block}{Correct Answer: B}
			SOC 2 Type II reports on control effectiveness over a period—the strongest assurance for security controls.
		\end{block}
	\end{frame}
	
	\begin{frame}{D4-Q615: Security Program Metrics}
		\begin{block}{Question}
			Which security program metric is MOST valuable for executives?
		\end{block}
		\begin{enumerate}[(A)]
			\item Number of firewall rules
			\item Trends in critical risks, incidents, and control effectiveness vs. tolerance
			\item Office location
			\item Marketing
		\end{enumerate}
	\end{frame}
	\begin{frame}{D4-Q615: Answer}
		\begin{block}{Correct Answer: B}
			Executives need trends in critical risks, incidents, and control effectiveness relative to tolerance—not operational counts.
		\end{block}
	\end{frame}

\end{document}